\documentclass[11pt]{article}
\usepackage[T1]{fontenc}
\usepackage{lmodern}
\usepackage[margin=1in]{geometry}
\usepackage{amsmath,amssymb,amsthm,mathtools}
\usepackage{aliascnt}
\usepackage{microtype}
\usepackage{booktabs}
\usepackage{tikz}
\usetikzlibrary{arrows.meta,calc,positioning}
\usepackage[numbers,sort&compress]{natbib}
\usepackage{xurl}
\usepackage[colorlinks=true,linkcolor=blue!45!black,citecolor=blue!45!black,urlcolor=blue!45!black]{hyperref}
\usepackage[capitalise,noabbrev]{cleveref}
\hypersetup{pdftitle={A radial continuum phase transition with algebraic decay},pdfauthor={OpenAI}}
\newtheorem{theorem}{Theorem}[section]
\newaliascnt{proposition}{theorem}
\newtheorem{proposition}[proposition]{Proposition}
\aliascntresetthe{proposition}
\newaliascnt{lemma}{theorem}
\newtheorem{lemma}[lemma]{Lemma}
\aliascntresetthe{lemma}
\newaliascnt{corollary}{theorem}
\newtheorem{corollary}[corollary]{Corollary}
\aliascntresetthe{corollary}
\theoremstyle{definition}
\newaliascnt{definition}{theorem}
\newtheorem{definition}[definition]{Definition}
\aliascntresetthe{definition}
\theoremstyle{remark}
\newaliascnt{remark}{theorem}
\newtheorem{remark}[remark]{Remark}
\aliascntresetthe{remark}
\newcommand{\R}{\mathbb R}
\newcommand{\N}{\mathbb N}
\newcommand{\dd}{\,\mathrm d}
\newcommand{\ee}{\mathrm e}
\DeclareMathOperator{\conv}{conv}
\DeclareMathOperator{\supp}{supp}
\DeclareMathOperator{\dist}{dist}
\DeclareMathOperator{\diam}{diam}
\DeclareMathOperator{\Lip}{Lip}
\newcommand{\norm}[1]{\left\lVert#1\right\rVert}
\newcommand{\abs}[1]{\left|#1\right|}
\newcommand{\ip}[2]{\left\langle#1,#2\right\rangle}
\newcommand{\clB}[2]{\overline B(#1,#2)}
\numberwithin{equation}{section}
\title{A radial continuum phase transition with algebraic decay}
\author{OpenAI}
\date{September 24, 2026}
\begin{document}
\maketitle
\begin{abstract}
We construct a bounded, continuous, stable radial pair potential in three
dimensions with \(|\phi(r)|\le Cr^{-3-1/32}\) for \(r\ge1\). Throughout an
open interval of positive densities, its canonical thermodynamic free energy
is finite at every positive inverse temperature and has a strict downward
derivative jump at one common finite positive inverse temperature.
\end{abstract}
\section{Introduction}\label{sec:introduction}

A phase transition is a singularity that appears in the thermodynamic
limit of a particle system.  We construct one for a radial pair
interaction with a bounded continuous core and an explicit integrable
power-law bound on its tail.

We consider identical classical particles in three-dimensional Euclidean space.
For a radial pair potential \(\phi\), a cube \(\Lambda_L=[0,L]^3\), and \(N\ge0\), set
\begin{equation}\label{eq:partition}
 U^\phi(x_1,\ldots,x_N)=\sum_{1\le i<j\le N}\phi(\abs{x_i-x_j}),
 \qquad
 Z^\phi_{L,N}(\beta)=\frac1{N!}
  \int_{\Lambda_L^N}\exp(-\beta U^\phi)\,\dd x_1\cdots\dd x_N,
\end{equation}
with \(Z^\phi_{L,0}=1\).
The position measure is fixed independently of the inverse temperature
\(\beta>0\).
A potential is \emph{stable} if \(U^\phi\ge-BN\) for a constant \(B\)
and every finite configuration.

Let \(p\) denote the packing constant
\[
 p=\lim_{L\to\infty}\frac{M_L}{L^3},
 \qquad
 M_L=\max\bigl\{\#X:X\subset[0,L]^3,\
               \abs{x-y}\ge1\ \text{for distinct }x,y\in X\bigr\}.
\]
The existence and positivity of \(p\) are proved in \Cref{f:packing};
its numerical value is not needed.

\begin{theorem}\label{thm:main}
There is a bounded continuous function
\(\phi:[0,\infty)\to\R\) and finite constants \(B,C>0\) such that
\[
 \sum_{i<j}\phi(\abs{x_i-x_j})\ge-BN,
 \qquad
 \abs{\phi(r)}\le C r^{-3-1/32}\quad(r\ge1).
\]
There are an open interval \(I\subset(0,\infty)\) centered at \(5p/3\)
and a common inverse temperature \(\beta_c\in[7/8,9/8]\) such that,
for every \(\rho\in I\), the canonical free energy
\[
 f(\beta,\rho)=
 -\lim_{L\to\infty}\frac1{\beta L^3}
       \log Z^\phi_{L,\lfloor\rho L^3\rfloor}(\beta)
\]
exists and is finite for every \(\beta>0\).
Its finite one-sided derivatives at \(\beta_c\) satisfy
\[
 f'_-(\beta_c,\rho)>f'_+(\beta_c,\rho).
\]
In particular, \(f(\,\cdot\,,\rho)\) is not differentiable at \(\beta_c\).
\end{theorem}

\subsection{Context and the role of the tail}

The existence of continuum phase transitions belongs to the program
highlighted by Simon's problems in mathematical physics
\citep{Simon1984}.  We address the explicit interaction class and
canonical singularity in \Cref{thm:main}.

Rigorous continuum coexistence results precede the present construction.
Ruelle proved a transition for a symmetric two-species classical system
\citep{Ruelle1971}.  Lebowitz, Mazel, and Presutti established a
liquid--vapor transition for a continuum model with an attractive
pair term and a repulsive four-body term
\citep{LebowitzMazelPresutti1998,LebowitzMazelPresutti1999}.
Dereudre and Renaud-Chan establish density coexistence for a class of
saturated continuum interactions
\citep[Theorems~1--2]{DereudreRenaudChan2026}.
Their radial-kernel example integrates over the union of particle
neighborhoods, giving a many-body energy
\citep[Section~2.3, Example~3]{DereudreRenaudChan2026}.
The species or many-body structures are part of these models; the
problem considered here uses a single species and a sum of radial pair
interactions.

Long-range weak attractions provide a classical route from microscopic
particle interactions to a nonconvex mean-field free energy.  The
Lebowitz--Penrose theory identifies the limiting free energy through
convexification of a reference free energy with a quadratic attraction
term \citep{LebowitzPenrose1966}.  Gerardi, Marchioro, Olivieri, and
Presutti extended this framework from hard-core references to
superstable interactions \citep{GerardiEtAl1973}.  In the present problem, the limit of
an interaction as its range tends to infinity is only an intermediate
object: the required singularity must belong to one fixed pair
potential.  Moreover, a pointwise power bound constrains how the masses
of its successive attractions may depend on their ranges.

The companion article \citep[Theorem~1.1]{OpenAICompanion2026} constructs
a radial potential with a divergent repulsive core, an integrable
\(o(r^{-3})\) tail, and a common transition inverse temperature for an
open interval of densities.  Both articles use small clusters and
separated supports of a pressure in two variables.  Here we prove the
bounded-core construction and the quantitative range estimate needed
for a fixed power margin.  All proofs required for the present theorem
are included below.  The two articles concern different potentials,
each with its own density interval and common critical inverse temperature.

A recent theorem of He, Jauslin, Lebowitz, and Peled proves a
liquid--vapor transition for a finite-range modified Kac pair
interaction \citep[Equation~(1.13), Theorems~1.3 and~1.5]{HeJauslinLebowitzPeled2026}.
Their interaction is defined using a fixed partition into boxes, and
the coexistence theorem is formulated at a fixed inverse temperature.
Our physical interaction is a radial function of particle separation;
boxes enter only the estimates used to analyze its integrals.  The
conclusion here is a strict inverse-temperature derivative jump of the
canonical free energy at one common critical point throughout an open
interval of densities.  It is an existence result for
the constructed infinite-range interaction, with no identification of
its phases as conventional fluid phases.

The estimates also have several classical ancestors.  The quadratic
occupation control is part of the superstability framework developed by Ruelle \citep{Ruelle1970}.
For our particular core, the needed bound follows from the weighted
graph maximization argument of Motzkin and Straus
\citep[Theorem~1]{MotzkinStraus1965}; we reproduce the argument and its
geometric application.  Supporting slopes and partial convex duality
are used in their standard finite-dimensional sense
\citep[Theorem~11.1 and Proposition~11.3]{RockafellarWets1998}.
The local occupation estimates, the quantitative attraction comparison,
and the persistence argument are proved at their precise strengths in
this paper.

\subsection{Proof strategy}

There are three tasks: construct a stable family of bounded potentials,
compare a long-range attraction with its ideal density transform, and
preserve a sufficiently large set of pressure supports while the
attraction ranges grow.  We describe how their outputs fit together.

\paragraph{A bounded reference family.}
A parameter \(t>1\) fixes the microscopic core; it is chosen once and is
independent of \(\beta\).  Below distance one, the core has height
\(t^2\) outside a narrow shell at distances comparable to \(\ee^{-t}\);
inside the shell it has a zero-energy plateau.  It vanishes beyond a
fixed radius.  Four particles can form a tetrahedron with all six pair
distances in the zero plateau, whereas five cannot lie in the shell.  This gives a quadratic lower bound
in the particle counts of small cubes.  The bound survives the addition
of nonnegative radial attractions with a prescribed total mass.
\Cref{f:section,f:local-section} use it to establish the canonical and
grand-canonical limits, uniform occupation moments, and continuity of
pressure under a small change of attraction mass.

For a potential \(\varphi\) in this family and a density \(d\ge0\),
write
\[
 H_\varphi(\beta,d)
 =-\lim_{L\to\infty}\frac{1}{tL^3}
        \log Z^\varphi_{L,\lfloor dL^3\rfloor}(\beta),
 \qquad
 Q_\varphi(\beta,h)
 =\sup_{d\ge0}\{hd-H_\varphi(\beta,d)\}.
\]
Thus \(H_\varphi=\beta f_\varphi/t\), and \(h\) is the field dual to
density.  The pressure \(Q_\varphi\) is convex jointly in
\((\beta,h)\).  A supporting slope \((e,d)\) at \((\beta_0,h_0)\)
means that
\[
 Q_\varphi(\beta,h)\ge Q_\varphi(\beta_0,h_0)
             +e(\beta-\beta_0)+d(h-h_0)
\]
for every \(\beta>0\) and \(h\in\R\).
Two such slopes with the same density \(d=\rho\) and distinct first
coordinates give two distinct upper supporting slopes of
\(H_\varphi(\,\cdot\,,\rho)\).  This is the route to the derivative
jump, proved explicitly in \Cref{m:canonical-corner}.

\paragraph{An attraction with controlled error.}
For a pressure \(Q\), an ideal attraction of mass \(a>0\) acts by
\begin{equation}\label{eq:intro-transform}
 (T_aQ)(\beta,h)
 =\sup_{u\in\R}
   \left\{Q(\beta,h+\beta au)-\frac{\beta a}{2}u^2\right\}.
\end{equation}
Square completion identifies its maximizing shifts \(u\) with
maximizing densities after addition of the quadratic attraction term.
\Cref{k:section-transform} proves attainment, uniform bounds on these
shifts, and an exact formula for all supporting slopes of \(T_aQ\).

Fix a nonnegative radial \(C^1\) kernel \(w\), supported in the unit
ball and of integral one, and put \(w_R(x)=R^{-3}w(x/R)\).
If the reference interaction \(\varphi_0\) has range at most \(R_0\),
then \Cref{k:range} compares the physical attraction \(-ta w_R\)
with the ideal transform:
\[
 \bigl\|Q_{\varphi_0-ta w_R}-T_aQ_{\varphi_0}\bigr\|_{\mathcal D}
 \le C_{\mathcal D,t}\sqrt{R_0/R}
\]
on each compact parameter set \(\mathcal D\).  The constant is
uniform over earlier reference interactions and over positive masses
\(a\) tending to zero.  Boxes of side comparable to
\(\sqrt{R_0R}\) balance the cost of deleting reference interactions
across boxes against the variation of the new kernel inside them.
For the lower bound, independent reference grand ensembles are placed
in these boxes.  Completing a square avoids any need to identify their
mean density with the maximizing shift, so the comparison remains valid
at coexistence without selecting a phase.

\paragraph{Three slope groups and one fixed potential.}
As \(t\) grows, the bare-core pressure tends to
\(\max\{0,ph,4ph-9p\}\).  Its branches come from vacuum, single
particles, and four-particle clusters, respectively.  A first ideal
attraction turns these into three affine branches with noncollinear
slopes in the \((e,d)\) plane.  \Cref{i:section} proves that, for one
large finite \(t\), the actual generating supports remain near these
three slopes throughout a parameter neighborhood.

Uniform convergence of convex pressures by itself need not preserve
a corner.  The additional control comes from the following property
of a maximizing shift in \eqref{eq:intro-transform}: at its shifted
parameter \((\beta,h+\beta au)\), the reference pressure has a
parabolic upper bound in the field variable
(\Cref{k:active-parabola}).  In \Cref{m:one-step} this bound confines
nearby supports to one density band and hence to one of the separated
slope groups.  An error \(\delta\) in pressure then costs
\(O(\delta/a)\) in the slopes.

After the first attraction, take \(R_{n+1}=R_n^2\) and
\(a_n=R_n^{-1/16}\).  The range errors are \(O(R_n^{-1/4})\), and
their ratios to the next masses are \(O(R_n^{-1/8})\).  All relevant
losses are summable.  \Cref{sec:iteration} consequently preserves a
common two-dimensional ball of pressure supports, passes it to the
limit, and uses horizontal slices near density \(5p/3\) to obtain the
canonical singularity at a common inverse temperature throughout a
density interval.  The same schedule gives
\(\sum_n a_nR_n^{1/32}<\infty\), which supplies the pointwise tail
bound.  The core, all attraction parameters, and the density interval
are fixed before the limiting critical parameter is selected.

\section{A stable family and its thermodynamic functions}
\label{f:section}

We first construct a bounded repulsive core which permits clusters of at
most four particles at very short separations.  Its quadratic control of
local particle numbers will be uniform over all attraction ranges used
later.  Throughout the proof, all position integrals use Lebesgue measure
in one fixed microscopic length unit, and all finite volumes have free
boundaries.

\subsection{Packing and the microscopic core}

\begin{lemma}[Packing density]\label{f:packing}
Let $M_L$ be the maximum cardinality of a subset of $[0,L]^3$ whose
distinct points have mutual distances at least one.  Then
\[
  p=\lim_{L\to\infty}\frac{M_L}{L^3}
\]
exists and belongs to $(0,\infty)$.
\end{lemma}

\begin{proof}
Disjoint balls of radius $1/3$ about the points lie in
$[-1/3,L+1/3]^3$, so $M_L\le C(L+1)^3$.  This also bounds the
possible cardinalities; compactness of the configuration space for each
cardinality shows that the maximum is attained.  The integer lattice
gives a positive lower asymptotic density.

Fix $s>0$ and repeat a maximizing configuration in $[0,s]^3$ by
translations in $(s+1)\mathbb Z^3$.  Points in distinct translates are
separated by at least one: in a coordinate in which the translates
differ, their coordinate separation is at least $(s+1)-s=1$.
Taking all complete translates contained in $[0,L]^3$ gives
\[
 \liminf_{L\to\infty}\frac{M_L}{L^3}
 \ge \frac{M_s}{(s+1)^3}.
\]
Letting $s\to\infty$ along a sequence realizing the upper asymptotic
density makes the right side tend to that upper density.  The lower
and upper limits therefore agree.
\end{proof}

Set $A=3/(2p)$ and $\Delta=10^{-3}$.  The parameter $t>1$ below is a
construction parameter, distinct from inverse temperature $\beta$.
We shall eventually choose one sufficiently large $t$ and keep it
fixed.  Put
\[
 l_t=\ee^{-t},\qquad I_t=[l_t,(1+\Delta)l_t],
 \qquad 2(1+\Delta)l_t<1.
\]
Define $\psi_t:[0,\infty)\to[0,t^2]$ by
\begin{equation}\label{f:core}
 \begin{aligned}
 \psi_t(r)&=t^2
   &&\text{on }[0,l_t]\cup[(1+\Delta)l_t,1],\\
 \psi_t(r)&=0
   &&\text{on }[(1+\Delta/3)l_t,(1+2\Delta/3)l_t]
                    \cup[1+1/t,\infty),
 \end{aligned}
\end{equation}
and make it affine on each of the three intervening intervals.  Thus
$\psi_t$ is continuous, including at zero, nonnegative, and supported
in $[0,2]$.

Fix a radial nonnegative $C^1$ function $w$ on $\R^3$, supported in the
closed unit ball, with $\int_{\R^3}w(x)\,\dd x=1$.  For example, one
may normalize $((1-\abs{x}^2)_+)^2$.  Write
$w_R(x)=R^{-3}w(x/R)$.  When its argument is a displacement in $\R^3$,
we also write $\psi_t$ for the radial lift $x\mapsto\psi_t(\abs{x})$.
The family of potentials used below is
\begin{equation}\label{f:family}
 \varphi(x)=\psi_t(\abs{x})-S(x),\qquad
 S(x)=t\sum_j\alpha_j w_{R_j}(x),\qquad
 R_j\ge1,\quad \alpha_j\ge0,\quad \sum_j\alpha_j\le A+1.
\end{equation}
The index set may be empty, finite, or countable.  Since
$\norm{w_{R_j}}_\infty\le\norm{w}_\infty$, the sum converges
uniformly in the countable case.  Every such potential is therefore
bounded, continuous, radial, and nonpositive when $\abs{x}\ge2$.

The following lemma combines an elementary geometric obstruction with
an extremal graph estimate.  The weighted edge maximization in its
proof is the Motzkin--Straus argument for Tur\'an's bound
\citep[Theorem~1]{MotzkinStraus1965}.  We include it to make the
quadratic core estimate explicit.

\begin{lemma}[The shell graph]\label{f:shell-graph}
For any finite configuration in $\R^3$, join two particles when their
distance belongs to $I_t$.  This graph has no clique of size five.
On any $m$ vertices it has at most $3m^2/8$ edges.
\end{lemma}

\begin{proof}
If five points formed a clique, rescale their distances by $l_t^{-1}$
and take the four difference vectors from one of them.  Their Gram
matrix differs entrywise by at most $2\Delta+\Delta^2$ from the matrix
with diagonal one and off-diagonal entries $1/2$.  For a diagonal entry
this follows from $1\le r^2\le(1+\Delta)^2$; for an off-diagonal entry
use polarization,
$\ip{v_i}{v_j}=(\abs{v_i}^2+\abs{v_j}^2-\abs{v_i-v_j}^2)/2$.
The reference matrix has smallest eigenvalue $1/2$.  The perturbation
has operator norm at most $4(2\Delta+\Delta^2)<1/2$, so the Gram matrix
would be positive definite and have rank four.  Four vectors in
$\R^3$ cannot have this Gram matrix.

For the edge bound, maximize $\sum_{\{i,j\}\in E}u_i u_j$ over
$u_i\ge0$, $\sum_i u_i=1$, choosing a maximizer with the smallest
possible support.  If two nonadjacent vertices have positive weights,
the objective is affine in their two weights with their sum fixed.
Putting their combined weight on one of the vertices cannot decrease
the objective, contradicting minimality of the support.  The support
is therefore a clique of some size $r\le4$.  Its objective is
\[
 \frac12\left(1-\sum_i u_i^2\right)
 \le\frac12(1-1/r)\le\frac38.
\]
For $m>0$, apply this bound to equal weights $1/m$.  The empty graph
case is immediate.
\end{proof}

\subsection{Quadratic control of cell counts}

We now convert the exclusion of five-particle shell cliques into a
lower bound that controls arbitrary particle numbers, including
coincident particles.  This is a concrete superstable estimate in the
sense of \citet{Ruelle1970}; the constants will be uniform over all
attraction ranges.

Partition space into half-open cubes $C_i$ of side $b=1/4$, with any
fixed translation of the grid.  When working in a bounded region,
replace each cell by its intersection with that region and omit empty
cells.  Let $n_i$ be the number of particles in cell $i$.  If
\[
 J_{ij}=\sup_{x\in C_i,\,y\in C_j}S(x-y),
\]
then $J$ is symmetric and
\begin{equation}\label{f:row-bound}
 \sup_i\sum_jJ_{ij}\le M:=c_b t(A+1),
 \qquad c_b=\norm{w}_\infty(2/b+5)^3.
\end{equation}
Indeed, a cell interacts through $w_R$ with at most $(2R/b+5)^3$
cells, each contributing at most $R^{-3}\norm{w}_\infty$.  Since
$R\ge1$, their row sum is at most $c_b$.  Summing the coefficients
proves \eqref{f:row-bound}, also for a countable sum and cropped cells.

We shall need the same estimates after removing some interactions.
More precisely, in any bounded region allow symmetric measurable
pair interactions
\begin{equation}\label{f:modified-energy}
 v(x,y)=c(x,y)\psi_t(\abs{x-y})-K(x,y),
 \qquad 0\le c\le1,\quad K\ge0,
\end{equation}
where $c(x,y)=1$ for points in the same small cell, $K$ is bounded,
and the symmetric matrix
$J_{ij}=\sup_{x\in C_i,y\in C_j}K(x,y)$ has row sums at most $M$.
The coefficients are fixed functions of the positions, independent of
the particle configuration.  This includes partial deletion of any
attractive pairs, partial deletion of core pairs between distinct
cells, and their linear interpolations whenever they share the stated
row bound.  No translation invariance is required for these finite
volume estimates.  Write $U^v(X)=\sum_{r<s}v(x_r,x_s)$, and write
$U^\varphi$ for a member of \eqref{f:family}.

\begin{proposition}[Common stability]\label{f:stability}
Set
\begin{equation}\label{f:stability-constants}
 D_0=t^2/8,\qquad B_0=t^2/2.
\end{equation}
For every configuration, including coincident particles,
\begin{equation}\label{f:cell-stability}
 U^{\psi_t}\ge D_0\sum_i n_i^2-B_0N,
 \qquad
 U^v\ge(D_0-M/2)\sum_i n_i^2-B_0N.
\end{equation}
Choose $t$ sufficiently large that
\begin{equation}\label{f:t-requirements}
 D_0\ge8M,\qquad ct>4(A+1),\qquad c=1/3456.
\end{equation}
Then, in any axis-parallel cube of side $L\ge1$, with $V=L^3$,
\begin{equation}\label{f:density-stability}
 U^v\ge ct^2\frac{N^2}{V}-\frac{t^2}{2}N.
\end{equation}
In particular all members of \eqref{f:family} are stable with the
same constant $B_0$.
\end{proposition}

\begin{proof}
The diameter $\sqrt3 b$ of a small cell is less than one.  By
\Cref{f:shell-graph}, among $m$ particles in one cell at least
$m(m-1)/2-3m^2/8=m^2/8-m/2$ pairs have distance outside $I_t$.
Each of these pairs contributes $t^2$, while all remaining core
contributions are nonnegative.  Summing over cells proves the first
bound, and also proves it when core terms between cells are deleted.
For the attractive contribution, including extra diagonal self terms
only enlarges the bound, and symmetry gives
\[
 \sum_{r<s}K(x_r,x_s)
 \le\frac12\sum_{i,j}J_{ij}n_i n_j
 \le\frac14\sum_{i,j}J_{ij}(n_i^2+n_j^2)
 \le\frac M2\sum_i n_i^2.
\]
This proves \eqref{f:cell-stability}.

There are at most $(L/b+2)^3\le216V$ nonempty cells.  Hence
$\sum_i n_i^2\ge N^2/(216V)$.  Under \eqref{f:t-requirements},
$D_0-M/2\ge t^2/16$, giving \eqref{f:density-stability} with the stated
$c$.  The requirements on $t$ are compatible: $M$ grows linearly in
$t$, whereas $D_0$ grows quadratically.  The constants $c$ and $c_b$
are independent of $t$ and of the attraction ranges.
\end{proof}

The two margins in \eqref{f:t-requirements} will be used separately.
The inequality $D_0\ge8M$ makes the influence of neighboring cell counts
small compared with the repulsion within one cell; this will yield
volume-independent occupation moments.  The inequality $ct>4(A+1)$
keeps the density bound coercive even after adding the quadratic
attractive weight in the next grand-sum estimate.

\subsection{Canonical limits and pressure duality}

The local quadratic bound controls configurations of high density.
We next use it to establish the thermodynamic functions required by
the attraction transform.  The proof applies directly to the infinite
attraction sum: separated boxes interact only attractively, so no
finite-range approximation is needed to construct these limits.

For $\Lambda_L=[0,L]^3$, define the canonical partition functions
and their normalization by
\begin{equation}\label{f:canonical-definition}
 \begin{aligned}
 Z^\varphi_{L,n}(\beta)
 &=\frac1{n!}\int_{\Lambda_L^n}\exp(-\beta U^\varphi(X))\,\dd X,
 & Z^\varphi_{L,0}&=1,\\
 H^\varphi_{L,n}(\beta)&=-\frac1{tL^3}\log Z^\varphi_{L,n}(\beta).
 \end{aligned}
\end{equation}

\begin{lemma}[Finite volume canonical bounds]\label{f:canonical-bounds}
There is an absolute constant $C<\infty$, independent of $t$ and all
attraction ranges, such that for every admissible $t$, $L\ge1$,
$n\ge0$, $\beta>0$, and $d_L=n/L^3$,
\begin{equation}\label{f:finite-bounds}
 \beta t(cd_L^2-d_L/2)+\frac{d_L\log d_L-d_L}{t}
 \le H^\varphi_{L,n}(\beta)
 \le \beta Ct d_L^2+\frac{d_L\log d_L}{t},
\end{equation}
where $0\log0=0$.
\end{lemma}

\begin{proof}
The lower bound follows by bounding the integrand from above using
\eqref{f:density-stability} and using $n!\ge(n/\ee)^n$ for $n\ge1$.
For the other bound, let $X_1,\ldots,X_n$ be independent uniform
points in $\Lambda_L$.  Since $U^\varphi\le U^{\psi_t}$ and
$\int_{\R^3}\psi_t(\abs{x})\,\dd x\le t^2\abs{B(0,2)}$,
\[
 \mathbb E U^\varphi
 \le\frac{n(n-1)}{2L^3}\int_{\R^3}\psi_t(\abs{x})\,\dd x
 \le Ct^2n^2/L^3.
\]
Jensen's inequality gives
$\log Z^\varphi_{L,n}\ge n\log L^3-\log n!-\beta Ct^2n^2/L^3$.
Now use $n!\le n^n$.  At $n=0$ both bounds are equalities.
\end{proof}

The following bound controls the entire large-particle-number tail.
The additional quadratic factor is included because it will also
occur in the range comparison.

\begin{lemma}[Grand sum tails]\label{f:grand-tail}
Fix an admissible $t$ and a compact $\mathcal D\subset\Omega$, where
$\Omega=(0,\infty)\times\R$ has coordinates $(\beta,h)$.
There are constants $D\ge1$ and $a_0>0$, depending only on $t$ and
$\mathcal D$, such that, for every $\varphi$ in \eqref{f:family},
$L\ge1$, and $0\le\zeta\le2(A+1)$,
\begin{equation}\label{f:grand-tail-bound}
 \sum_{n>DL^3}\ee^{thn}Z^\varphi_{L,n}(\beta)
          \exp\left(\frac{\beta t\zeta n^2}{2L^3}\right)
 \le\frac{\exp[-a_0(\lfloor DL^3\rfloor+1)]}{1-\ee^{-a_0}}
 \qquad ((\beta,h)\in\mathcal D).
\end{equation}
In particular the tail tends uniformly to zero, and is at most one
for all sufficiently large $L$.  These constants require no local
occupation estimate.
\end{lemma}

\begin{proof}
Enclose $\mathcal D$ in $\beta\in[\beta_-,\beta_+]$ and
$\abs{h}\le H$.  Put $V=L^3$.  By \eqref{f:finite-bounds}, the logarithm
of the summand at $n\ge1$ is at most
\[
 -\beta t(ct-\zeta/2)\frac{n^2}{V}
 +(th+\beta t^2/2+1)n-n\log(n/V).
\]
The choice \eqref{f:t-requirements} implies
$ct-\zeta/2\ge ct/2$.  Set
$\kappa=\beta_-ct^2/2$ and $B=tH+\beta_+t^2/2+1$.
For $n/V\ge1$ the displayed expression is at most
$-\kappa n^2/V+Bn$.  Choose $D\ge\max(1,2B/\kappa)$.
For $n>DV$ this is at most
$-\kappa n^2/(2V)\le-(\kappa D/2)n$.
Summing this geometric majorant proves the assertion with
$a_0=\kappa D/2$.
\end{proof}

\begin{proposition}[Thermodynamic limits and duality]
\label{f:thermodynamics}
For every member $\varphi$ of \eqref{f:family}, including countably
infinite sums, and every $\beta>0$, $d\ge0$, the limit
\begin{equation}\label{f:canonical-limit}
 H_\varphi(\beta,d)
 =\lim_{\substack{L\to\infty\\ n/L^3\to d}}
      H^\varphi_{L,n}(\beta)
\end{equation}
exists along every such sequence of real sides and integer counts.
It is finite, convex and continuous in $d\in[0,\infty)$, and concave
in $\beta>0$ for each fixed $d$.  The bounds \eqref{f:finite-bounds}
hold with $H_\varphi(\beta,d)$ and $d$ in place of their finite volume
counterparts.

The grand partition function and pressure satisfy
\begin{equation}\label{f:pressure-definition}
 \Xi^\varphi_L(\beta,h)
 =\sum_{n\ge0}\ee^{thn}Z^\varphi_{L,n}(\beta),
 \qquad
 Q_\varphi(\beta,h)
 =\lim_{L\to\infty}\frac{\log\Xi^\varphi_L(\beta,h)}{tL^3}.
\end{equation}
The pressure is finite, nonnegative, and jointly convex on $\Omega$.
For $\beta>0$,
\begin{equation}\label{f:duality}
 \begin{split}
 Q_\varphi(\beta,h)
   &=\sup_{d\ge0}\{hd-H_\varphi(\beta,d)\},\\
 H_\varphi(\beta,d)
   &=\sup_{h\in\R}\{hd-Q_\varphi(\beta,h)\},\qquad d>0.
 \end{split}
\end{equation}
For fixed admissible $t$, pressures of the whole family are uniformly
bounded on each compact subset of $\Omega$.
\end{proposition}

\begin{proof}
We first prove \eqref{f:canonical-limit}, without using a pressure
limit or an approximation by finite-range potentials.  At $d=0$,
both bounds in \eqref{f:finite-bounds} tend to zero.  Fix $d>0$ and
$\beta>0$.  Let $H_*$ be the lower limit of
$H^\varphi_{L,n}(\beta)$ as $L\to\infty$ and $n/L^3\to d$, and choose
$s_j\to\infty$ and integers $n_j$ with
\[
 n_j/s_j^3\longrightarrow d,\qquad
 H^\varphi_{s_j,n_j}(\beta)\longrightarrow H_*.
\]
The finite volume bounds make $H_*$ finite.  Consider any other
sequence $L\to\infty$, $n/L^3\to d$.  For each fixed, sufficiently
large $j$, place
$k=\lfloor L/(s_j+2)\rfloor^3$ cubes of side $s_j$ in $\Lambda_L$,
separated by gaps of width at least two.  For any nonnegative integers
$m_1,\ldots,m_k$ summing to $n$, restriction to configurations with
these box counts gives
\begin{equation}\label{f:canonical-product}
 Z^\varphi_{L,n}(\beta)
 \ge\prod_{i=1}^k Z^\varphi_{s_j,m_i}(\beta).
\end{equation}
Indeed, interactions between different boxes are nonpositive.  There
are $n!/\prod_i m_i!$ disjoint assignments of particle labels to the
boxes, so the global factor $1/n!$ leaves exactly the product of
their factors $1/m_i!$.

We retain the favorable count $n_j$ in most boxes and correct the
total particle number in the remaining boxes.  Only a vanishing fraction
of boxes will change, and \eqref{f:finite-bounds} will bound their
canonical costs uniformly.  To make this exact for every target count, put
\[
 q_j=\lfloor(4d+4)s_j^3\rfloor,
 \qquad c_d=\min(d/2,1).
\]
For all sufficiently large $j$,
$n_j\ge c_ds_j^3$ and $q_j-n_j\ge c_ds_j^3$.
Start with $n_j$ particles in every box.  If
$E=n-kn_j\ge0$, fill boxes successively to $q_j$, with at most one
last box given an intermediate integer count.  If $E<0$, empty boxes
successively, again with at most one last intermediate count.  This
changes at most $r$ boxes, where
\begin{equation}\label{f:integer-adjustment}
 \frac rk\le\frac1{c_d}
       \abs{\frac{n}{ks_j^3}-\frac{n_j}{s_j^3}}+\frac1k.
\end{equation}
For fixed large $j$ and then sufficiently large $L$, the target lies
between $0$ and $kq_j$: its mean count divided by $s_j^3$ tends to
$d(1+2/s_j)^3<q_j/s_j^3$.  Thus the adjustment never exhausts the
available boxes.  The quotient and remainder in division of $E$ by
$q_j-n_j$, or of $-E$ by $n_j$, give the asserted integer counts.

All counts used have density in $[0,4d+4]$.  By
\eqref{f:finite-bounds}, their normalized canonical costs have absolute
value at most some $K<\infty$, uniformly for large $j$.  Taking
logarithms in \eqref{f:canonical-product} consequently yields
\[
 H^\varphi_{L,n}(\beta)
 \le\frac{ks_j^3}{L^3}
       \left(H^\varphi_{s_j,n_j}(\beta)+2K\frac rk\right).
\]
First let $L\to\infty$ with $j$ fixed.  The prefactor tends to
$(s_j/(s_j+2))^3$, and the upper limit of $r/k$ is at most
\[
 c_d^{-1}\abs{d(1+2/s_j)^3-n_j/s_j^3}.
\]
Then let $j\to\infty$.  This error tends to zero and the preceding
upper bound tends to $H_*$.  Every sequence has lower limit at least
$H_*$ by its definition; thus every sequence has the same limit.
This proves \eqref{f:canonical-limit} with the required order of limits.

For density convexity, fix $d_1,d_2\ge0$ and $\theta\in[0,1]$.
Choose $s_j\to\infty$, integers $m_{r,j}$ with
$m_{r,j}/s_j^3\to d_r$ for $r=1,2$, and integers $q_j\to\infty$.
In a cube of side $q_j(s_j+2)$, use $q_j^3$ separated boxes of side
$s_j$, assigning count $m_{1,j}$ to $\lfloor\theta q_j^3\rfloor$
of them and count $m_{2,j}$ to the rest.  The total density tends to
$\theta d_1+(1-\theta)d_2$.  The product inequality and the already
proved all-sequence limit, applied directly to this diagonal sequence,
give
\[
 H_\varphi(\beta,\theta d_1+(1-\theta)d_2)
 \le\theta H_\varphi(\beta,d_1)
       +(1-\theta)H_\varphi(\beta,d_2).
\]
This proof does not presuppose continuity in density.  A finite convex
function is continuous at positive densities: on any interior compact
interval its secant slopes are bounded by secants to two exterior
points.  Continuity at zero follows from \eqref{f:finite-bounds}.
H\"older's inequality makes $\log Z^\varphi_{L,n}(\beta)$ convex in
$\beta$.  Passing to the canonical limit proves the stated concavity.

The grand sums are finite by \Cref{f:grand-tail}, with $\zeta=0$.
Fix $(\beta,h)$.  Choose $D$ there and write $V=L^3$.  Above $DV$ the
entire sum tends to zero.  Among the remaining
$\lfloor DV\rfloor+1$ terms choose a maximizing count $n_L$; the
maximum is at least one because it includes the vacuum.  Therefore
\[
 \frac{\log\Xi^\varphi_L}{tV}
 \le h\frac{n_L}{V}-H^\varphi_{L,n_L}(\beta)
       +\frac{\log(\lfloor DV\rfloor+2)}{tV}
\]
for all sufficiently large $L$.  From any subsequence choose a further
subsequence with $n_L/V\to d\in[0,D]$.  The canonical limit bounds
its upper limit by $hd-H_\varphi(\beta,d)$.  Conversely, for every
fixed $d\ge0$ the single term $n=\lfloor dV\rfloor$ bounds the lower
limit by $hd-H_\varphi(\beta,d)$.  This proves existence of the
pressure and the first duality in \eqref{f:duality}.

The vacuum gives $Q_\varphi\ge0$.  The canonical lower bound and
$d\log d-d\ge-1$ bound the first supremum in \eqref{f:duality} by a
finite quadratic maximum, uniformly over the family on any fixed
parameter compact.  H\"older's inequality on the disjoint union of
all configuration spaces makes each $\log\Xi^\varphi_L$ jointly
convex in $(\beta,h)$; its pointwise limit has the same property.

Finally fix $d>0$.  Convexity of the finite function
$H_\varphi(\beta,\cdot)$ gives a supporting slope $h$ there.  Explicitly,
one may choose $h$ between
\[
 \sup_{0\le v<d}\frac{H_\varphi(\beta,d)-H_\varphi(\beta,v)}{d-v}
 \quad\text{and}\quad
 \inf_{v>d}\frac{H_\varphi(\beta,v)-H_\varphi(\beta,d)}{v-d}.
\]
The left endpoint is at most the right endpoint by convexity, and
both are finite by comparison with secants through fixed points on
either side of $d$.  The resulting supporting inequality holds for
all $v\ge0$, so the first duality gives
$Q_\varphi(\beta,h)=hd-H_\varphi(\beta,d)$.  For all other $h$ the
same first duality gives the opposite variational inequality.  This
proves the second equality in \eqref{f:duality} without a differentiability
or uniqueness assumption.
\end{proof}

\section{Occupation bounds and perturbation estimates}
\label{f:local-section}

The next estimates turn the quadratic stability bound into uniform
control of local particle numbers and changes of pressure.  They play
the role of the occupation bounds in the theory of superstable
interactions \citep{Ruelle1970}.  We prove them for the precise family
of deletions and interpolations used in the range comparison.

From this section onward the occupation estimates are used with $t$
fixed.  Their constants may depend on $t$ and on a compact set of
$(\beta,h)$, but never on the number or sizes of the attraction ranges.
The finite volume systems in \eqref{f:modified-energy} are included;
in particular, deleting interactions between larger boxes will not
require a new thermodynamic limit theorem.

If a bounded region meets $k$ small cells, its grand partition
function has the exact decomposition
\begin{equation}\label{f:cell-measure}
 \Xi^v(\beta,h)
 =\sum_{n_1,\ldots,n_k\ge0}
    \frac{\ee^{th\sum_i n_i}}{\prod_i n_i!}
    \int_{\prod_i C_i^{n_i}}\exp(-\beta U^v(X))\,\dd X.
\end{equation}
Indeed, the $N!/\prod_i n_i!$ assignments of $N=\sum_i n_i$ labels
to disjoint cells cancel the global $1/N!$.  This identity also gives
the conditional grand distribution within one cell after the outside
configuration has been specified.

\begin{proposition}[Uniform occupation moments]\label{f:moments}
Fix $t$ satisfying \eqref{f:t-requirements} and a compact
$\mathcal D\subset\Omega$.  There are constants
$C_{\exp},C_1<\infty$, depending only on $t$ and $\mathcal D$, such
that every grand canonical measure for an interaction
\eqref{f:modified-energy} satisfies
\begin{equation}\label{f:moment-estimates}
 \sup_i\mathbb E\exp(\beta D_0n_i^2/4)\le C_{\exp},
 \qquad \sup_{i,j}\mathbb E[n_i n_j]\le C_1
 \qquad ((\beta,h)\in\mathcal D).
\end{equation}
The bounds hold for every bounded region with finitely many cropped
cells, and uniformly for all permitted deletions and interpolations.
\end{proposition}

\begin{proof}
If the region meets no small cell, its grand measure is the vacuum and
the assertion is immediate.  Otherwise, let the number of cells be positive.
Write $\alpha=D_0-M/2$, $u_0=\beta D_0/4$, and enclose
$\mathcal D$ in $\beta\in[\beta_-,\beta_+]$, $\abs{h}\le H$.
We first check, before taking any maximum, that
$F_i:=\mathbb E\exp(u_0n_i^2)$ is finite in each fixed volume.
By \eqref{f:cell-stability}, insertion of this exponential into
\eqref{f:cell-measure} bounds its numerator by the product of $k$
one-variable series with exponents
\[
 -\beta(\alpha-D_0/4)n_i^2
 -\beta\alpha\sum_{j\ne i}n_j^2
 +(th+\beta B_0)\sum_j n_j.
\]
Each spatial integral contributes at most
$\prod_j(b^3)^{n_j}$, with the factors $1/n_j!$ already in
\eqref{f:cell-measure}.  The quadratic coefficients are positive in
magnitude because
\[
 \alpha-D_0/4=3D_0/4-M/2\ge11D_0/16>0.
\]
Thus all these series are finite, and the denominator $\Xi^v$ is at
least one.  In particular,
$F_{\max}:=\max_iF_i<\infty$.  To bound this maximum independently of
the number of cells, we shall prove
\[
 F_{\max}\le C_2 F_{\max}^{\eta_0},
 \qquad \eta_0=4M^2/D_0^2\le1/16<1,
\]
where $C_2$ depends only on $t$ and $\mathcal D$.
The ratio of attraction strength to core repulsion will therefore
make the dependence on neighboring moments sublinear.

Condition on the configuration $Y$ outside cell $i$.  For $n$
particles inside that cell, the part of the energy involving at least
one inside particle is bounded below by
\begin{equation}\label{f:conditional-energy}
 (D_0-J_{ii}/2)n^2-B_0n-nG_i(Y),
 \qquad G_i(Y)=\sum_{j\ne i}J_{ij}n_j.
\end{equation}
This uses the complete core estimate inside the cell, discards
nonnegative core terms across cells, and bounds all attractions by
their cell suprema.  Dividing the conditional weights by the
outside-only Boltzmann factor leaves a denominator containing the
empty-cell term one.  Hence the conditional exponential moment is
at most its unnormalized numerator.  The inequality
\[
 nG_i\le D_0n^2/4+G_i^2/D_0
\]
and the factor $\exp(u_0n^2)$ leave the quadratic coefficient
$D_0/2-J_{ii}/2\ge7D_0/16$.  Consequently
\begin{equation}\label{f:conditional-moment}
 \mathbb E[\exp(u_0n_i^2)\mid Y]
 \le C_2\exp(\beta G_i(Y)^2/D_0),
\end{equation}
where the finite constant
\[
 C_2=\sum_{n\ge0}\frac{b^{3n}}{n!}
       \exp\left(-\frac{7\beta_-D_0}{16}n^2
                    +(tH+\beta_+B_0)n\right)
\]
depends only on $t$ and $\mathcal D$ and is at least one.

To close this estimate, Cauchy--Schwarz with the weights $J_{ij}$
and the row bound give
\[
 G_i^2\le M\sum_{j\ne i}J_{ij}n_j^2.
\]
Set $\theta_j=J_{ij}/M$ for $j\ne i$ and
$\theta_\varnothing=1-\sum_{j\ne i}\theta_j\ge0$.
With the exponent $\eta_0$ chosen above,
$\beta G_i^2/D_0\le\eta_0u_0\sum_{j\ne i}\theta_jn_j^2$.
Apply convexity of the exponential with the additional zero argument
of weight $\theta_\varnothing$, and then the concavity of
$x\mapsto x^{\eta_0}$, to obtain
\[
 \begin{split}
 \mathbb E\exp(\beta G_i^2/D_0)
 &\le\theta_\varnothing+
       \sum_{j\ne i}\theta_j\mathbb E\exp(\eta_0u_0n_j^2)\\
 &\le\theta_\varnothing+\sum_{j\ne i}\theta_jF_j^{\eta_0}
 \le F_{\max}^{\eta_0},
 \end{split}
\]
where the last inequality uses $F_{\max}\ge1$.
Taking expectations in \eqref{f:conditional-moment} and then the
maximum over cells yields
$F_{\max}\le C_2F_{\max}^{\eta_0}$, and therefore
\[
 F_{\max}\le C_2^{1/(1-\eta_0)}=:C_{\exp}.
\]
This bound is independent of the number of cells.  Since
$\exp(u_0n_i^2)\ge u_0n_i^2$ and
$u_0\ge\beta_-D_0/4$, the second moments are uniformly bounded.
Cauchy--Schwarz gives the mixed moment bound in
\eqref{f:moment-estimates}.

Cropping a cell only decreases its volume and the relevant suprema.
Every deletion or interpolation allowed in \eqref{f:modified-energy}
retains the full intra-cell core and the same row bound.  All steps
above therefore apply without change to those systems.
\end{proof}

\begin{corollary}[Absolute perturbation bound]\label{f:perturbation}
In a fixed bounded region, let $v_\theta=(1-\theta)v_0+\theta v_1$,
$0\le\theta\le1$, be bounded symmetric pair interactions satisfying
the hypotheses of \Cref{f:moments} with the same $t$ and $M$.
Let $J^{\mathrm{diff}}$ be a symmetric nonnegative matrix such that
\[
 \abs{v_1(x,y)-v_0(x,y)}\le J^{\mathrm{diff}}_{ij}
 \quad (x\in C_i,\ y\in C_j).
\]
Then, uniformly for $(\beta,h)\in\mathcal D$,
\begin{equation}\label{f:interpolation-bound}
 \abs{\log\Xi^{v_1}(\beta,h)-\log\Xi^{v_0}(\beta,h)}
 \le\frac{\beta C_1}{2}\sum_{i,j}J^{\mathrm{diff}}_{ij}.
\end{equation}
Thus signed deletions are controlled by the absolute interaction
difference, with no cancellation assumption.
\end{corollary}

\begin{proof}
Differentiation in $\theta$ is justified directly in finite volume.
The empty region is immediate.  Put $\alpha=D_0-M/2>0$.
If there are $k\ge1$ cells and the region has volume $V$, the common
energy bound gives
$U^{v_\theta}\ge\alpha N^2/k-B_0N$, while boundedness of the pair
difference gives $\abs{\partial_\theta U^{v_\theta}}\le C N^2$.
Uniformly in $\theta$ and on $\mathcal D$, the differentiated grand
series is therefore dominated by a constant times
\begin{equation}\label{f:derivative-domination}
 \sum_{N\ge0}(1+N^2)\frac{V^N}{N!}
       \exp\left(-\beta_-\alpha N^2/k
                      +(tH+\beta_+B_0)N\right)<\infty.
\end{equation}
The volume and the constant in this domination are fixed only for
the differentiation; the resulting estimate below is uniform in
volume.  Hence
\[
 \frac{\dd}{\dd\theta}\log\Xi^{v_\theta}
       =-\beta\mathbb E_\theta\partial_\theta U^{v_\theta}.
\]
For every configuration,
\[
 \abs{\partial_\theta U^{v_\theta}}
 \le\frac12\sum_{i,j}J^{\mathrm{diff}}_{ij}n_i n_j,
\]
where the diagonal uses $n_i(n_i-1)\le n_i^2$.
Apply \Cref{f:moments} and integrate over $\theta\in[0,1]$.
\end{proof}

\begin{corollary}[Pressure control by attraction mass]
\label{f:tail-pressure}
For a finite or countably infinite member of \eqref{f:family}, let
\[
 \varphi_k=\psi_t-t\sum_{j\le k}\alpha_jw_{R_j},
 \qquad m_k=\sum_{j>k}\alpha_j.
\]
For every compact $\mathcal D\subset\Omega$ there is a constant
$C_{\mathcal D,t}<\infty$, independent of the ranges, their number,
and $k$, such that
\begin{equation}\label{f:tail-pressure-bound}
 0\le Q_\varphi(\beta,h)-Q_{\varphi_k}(\beta,h)
 \le C_{\mathcal D,t}m_k
 \qquad ((\beta,h)\in\mathcal D).
\end{equation}
The same upper bound holds for the difference of normalized finite
volume log grand partition functions in every cube of side $L\ge1$.
\end{corollary}

\begin{proof}
Interpolate the whole tail at once:
\[
 \varphi_{k,\theta}
 =\varphi_k-\theta t\sum_{j>k}\alpha_jw_{R_j},
 \qquad 0\le\theta\le1.
\]
Even for infinitely many terms, the tail converges uniformly and is
a bounded continuous pair interaction, of supremum norm at most
$t\norm{w}_\infty m_k$.  Thus the derivative exists for each finite
configuration and the domination \eqref{f:derivative-domination}
justifies differentiation directly, including the infinite endpoint.
All intermediate attractions obey the original mass bound.

By the proof of \eqref{f:row-bound}, the difference matrix can be
chosen with
\[
 \sup_i\sum_jJ^{\mathrm{diff}}_{ij}\le c_btm_k.
\]
There are at most $216L^3$ cells.  Dividing
\eqref{f:interpolation-bound} by $tL^3$ gives the claimed finite
volume upper bound, for instance with
$C_{\mathcal D,t}=108\beta_+C_1c_b$.  Its lower bound follows because
adding a nonnegative attraction decreases the energy.

\Cref{f:thermodynamics} already supplies thermodynamic pressure
limits for both endpoints, including the infinite sum.  Passing to
those limits proves \eqref{f:tail-pressure-bound}; convergence of
truncation pressures is a consequence, not an input, of this estimate.
\end{proof}

\section{Ideal attraction and supporting slopes}
\label{k:section-transform}

An attraction spread over a very large region adds a quadratic density term
to the variational formula for pressure.  We first study that variational
operation, retaining the full subdifferential at points where the pressure
need not be differentiable.  Throughout this section, $Q=Q_\varphi$ and
$H=H_\varphi$ for a member of \eqref{f:family}.

\begin{definition}[Ideal attraction]
For $a>0$, define
\begin{equation}
 T_aQ(\beta,h)
 =\sup_{u\in\R}\left\{
 Q(\beta,h+\beta au)-\frac{\beta a}{2}u^2\right\}.
 \label{k:transform}
\end{equation}
An \emph{active shift} at $\lambda=(\beta,h)$ is a value of $u$ attaining
this supremum.  Its shifted parameter is
$x_u=(\beta,h+\beta au)$.
\end{definition}

\begin{proposition}[Attainment and uniform bounds]
\label{k:transform-properties}
Suppose the size conditions on $t$ in \Cref{f:stability} hold, including
$ct>4(A+1)$.  For every $0<a\le 2(A+1)$,
\begin{equation}
 T_aQ(\beta,h)
 =\max_{d\ge0}
 \left\{hd-H(\beta,d)+\frac{\beta a}{2}d^2\right\}.
 \label{k:density-transform}
\end{equation}
The active shifts in \eqref{k:transform} are exactly the maximizing
densities in \eqref{k:density-transform}; in particular, they are
nonnegative and form a nonempty compact set.  On each compact
$\mathcal D\subset\Omega$, they lie in a common interval $[0,U]$,
independently of $\varphi$ and $a$.  The constant $U$ may also be chosen
uniformly for all sufficiently large $t$.

The function $T_aQ$ is finite and convex on $\Omega$, and
\begin{equation}
 0\le T_aQ-Q\le C_{\mathcal D}a
 \qquad\hbox{on }\mathcal D.
 \label{k:small-mass}
\end{equation}
Here $C_{\mathcal D}$ can be chosen uniformly for sufficiently large $t$.
For fixed $t$, the pressures $Q$ and transforms $T_aQ$ have common bounds
on compact subsets of $\Omega$, and common local Lipschitz bounds,
uniformly over the same family and masses.
\end{proposition}

\begin{proof}
The pressure duality in \Cref{f:thermodynamics} gives, for every $u\in\R$,
\begin{align}
 Q(\beta,h+\beta au)-\frac{\beta a}{2}u^2
 &=\sup_{d\ge0}\left\{
 hd-H(\beta,d)+\frac{\beta a}{2}d^2
       -\frac{\beta a}{2}(u-d)^2\right\}.
 \label{k:square-completion}
\end{align}
Taking the supremum in $u$ proves \eqref{k:density-transform} initially
with a supremum.  To prove attainment and uniformity, use the lower bound
of \Cref{f:canonical-bounds} in the thermodynamic limit.  Since
$-d\log d+d\le1$ for $d\ge0$, its density objective is at most
\begin{align}
 hd-H(\beta,d)+\frac{\beta a}{2}d^2
 &\le-\beta\left(ct-\frac a2\right)d^2
       +\left(\abs h+\frac{\beta t}{2}\right)d+\frac1t\notag\\
 &\le-\frac{\beta ct}{2}d^2
       +\left(\abs h+\frac{\beta t}{2}\right)d+\frac1t.
 \label{k:coercivity}
\end{align}
This upper bound tends uniformly to $-\infty$ as $d\to\infty$ on every
parameter compact, whereas the objective at $d=0$ is zero.  Continuity
of $H(\beta,\cdot)$, including at zero, proves attainment and a common
bound on all maximizers.  Dividing \eqref{k:coercivity} by $\beta t$
gives the upper bound
\[
 -\frac c2d^2+
 \left(\frac{\abs h}{\beta t}+\frac12\right)d+
 \frac1{\beta t^2}.
\]
Its coefficients are uniformly bounded for $t\ge t_0>1$ and $\lambda$
in a fixed compact.  Thus the maximizing-density bound is also uniform
in sufficiently large $t$.

If $d$ maximizes \eqref{k:density-transform}, setting $u=d$ in
\eqref{k:square-completion} makes $u$ active.  Conversely, fix an active
$u$.  The density supremum defining $Q(\beta,h+\beta au)$ is attained:
its fixed shifted field and the canonical quadratic bound give coercivity.
For an attaining density $d$, \eqref{k:square-completion} can equal
$T_aQ(\beta,h)$ only if $d$ maximizes \eqref{k:density-transform} and
$(u-d)^2=0$.  This proves the claimed exact identity of maximizing sets
and also attainment in \eqref{k:transform}.

For fixed $u$, the branch
\[
 F_u(\beta,h)=Q(\beta,h+\beta au)-\beta au^2/2
\]
is convex: its parameter map is invertible and linear, and its last
term is linear.  The supremum of these branches is convex and is finite
by \eqref{k:coercivity}.  If $d$ is a maximizing density, pressure
duality and the nonnegativity of the added quadratic give
\[
 Q(\beta,h)\le T_aQ(\beta,h)
 \le Q(\beta,h)+\frac{\beta a}{2}d^2.
\]
This proves \eqref{k:small-mass}, with
$C_{\mathcal D}=\tfrac12\sup_{\mathcal D}\beta\, U^2$.
The canonical bound and pressure duality bound $Q$ uniformly on every
compact at fixed $t$; \eqref{k:small-mass} does the same for $T_aQ$.
Convexity then gives the asserted local Lipschitz bounds.
\end{proof}

We recall the elementary convex facts used below.  For a finite convex
function $F$ on $\Omega$, write
\[
 \partial F(\lambda)=
 \{v\in\R^2:F(y)\ge F(\lambda)+\ip{v}{y-\lambda}
                         \text{ for all }y\in\Omega\}.
\]
Such a function is continuous and locally Lipschitz.  Indeed, convexity
first bounds it above on a simplex by its values at the vertices;
reflection about an interior point supplies a lower bound on a smaller
neighborhood, and one-dimensional secant inequalities give a Lipschitz
bound on a still smaller neighborhood.  The directional derivative
\[
 F'(\lambda;z)=\lim_{r\downarrow0}
       \frac{F(\lambda+rz)-F(\lambda)}r
\]
is finite and sublinear in $z$.  A linear functional dominated by it
exists by linear separation (equivalently, the finite-dimensional
Hahn--Banach theorem).  The secant inequalities show that its vector is
a member of $\partial F(\lambda)$.  Thus this set is nonempty.  It is
convex and closed by its defining inequalities, and locally bounded by
the local Lipschitz bound, hence compact.  Those same inequalities and
continuity show that its graph is closed.  In particular, uniform bounds
for convex functions on a larger compact neighborhood give uniform bounds
for all their supporting slopes on a smaller compact neighborhood.
We write a pressure support as $(e,d)$, with coordinates dual to
$(\beta,h)$: $d$ is the density coordinate and $e$ the
inverse-temperature slope coordinate.

\begin{lemma}[Parabolic bound at an active point]
\label{k:active-parabola}
Let $a>0$ be as in \Cref{k:transform-properties}, and let $u$ be active
at $\lambda=(\beta,h)$.  Then for every $r\in\R$,
\begin{equation}
 Q(x_u+(0,r))\le Q(x_u)+ur+\frac{r^2}{2\beta a}.
 \label{k:parabola}
\end{equation}
Every member of $\partial Q(x_u)$ has second coordinate $u$.
\end{lemma}

\begin{proof}
Compare the active shift $u$ with $u+r/(\beta a)$ in
\eqref{k:transform}.  Expanding the two quadratic penalties proves
\eqref{k:parabola}.  If $(e,d)\in\partial Q(x_u)$, the supporting
inequality gives
\[
 dr\le Q(x_u+(0,r))-Q(x_u)
       \le ur+\frac{r^2}{2\beta a}.
\]
Divide separately for $r>0$ and $r<0$ and let $r\to0$ to obtain $d=u$.
\end{proof}

The subdifferential of a supremum of convex functions is governed by
its active branches.  General forms of this principle appear in
\citet[p.~864]{HantouteLopezZalinescu2008}.  Here the compactness of
active shifts and the parabolic bound give the following exact formula;
we prove both inclusions for this setting, including nonsmooth pressures.

\begin{proposition}[Exact subdifferential formula]
\label{k:subdifferential}
For each fixed $0<a\le2(A+1)$ and $\lambda=(\beta,h)\in\Omega$,
\begin{equation}
 \partial(T_aQ)(\lambda)
 =\conv\!\bigcup_{u\text{ active at }\lambda}
 \left\{\left(e+\frac a2u^2,u\right):(e,u)\in\partial Q(x_u)\right\}.
 \label{k:support-hull}
\end{equation}
The union on the right is compact, and every one of its members is an
actual supporting slope of $T_aQ$ at $\lambda$.
\end{proposition}

\begin{proof}
For fixed $u$, put $L_u(\beta,h)=(\beta,h+\beta au)$.  The defining
support inequalities, applied to $L_u$ and its inverse, give the exact
formula
\[
 \partial F_u(\lambda)=
 \left\{(e+aud-au^2/2,d):(e,d)\in\partial Q(L_u\lambda)\right\}.
\]
If $u$ is active, \Cref{k:active-parabola} sets $d=u$, giving the
generators in \eqref{k:support-hull}.  Since $F_u\le T_aQ$ everywhere
and $F_u(\lambda)=T_aQ(\lambda)$, each such generator supports $T_aQ$.
The convex hull of the displayed union is therefore contained in the
left side.

Denote that union by $G(\lambda)$.  The active set is bounded by
\Cref{k:transform-properties} and closed by continuity of $F_u(\lambda)$
in $u$, so it is compact.  Its shifted points form a compact subset of
$\Omega$.  The boundedness and closed graph of $\partial Q$ on this set
show that $G(\lambda)$ is compact.  Its convex hull is compact as well:
in $\R^2$, each point of a convex hull is a convex combination of at most
three points, by eliminating affine dependencies from any larger finite
combination.

For the reverse inclusion, fix $v_0\in\partial(T_aQ)(\lambda)$ and a
direction $z\in\R^2$.  For small $r>0$, let $\lambda_r=\lambda+rz$,
choose an active shift $u_r$ there, and choose
$v_r\in\partial F_{u_r}(\lambda_r)$.  The branch support at the endpoint
gives
\begin{equation}
 \ip{v_0}{z}
 \le\frac{T_aQ(\lambda_r)-T_aQ(\lambda)}r
 \le\frac{F_{u_r}(\lambda_r)-F_{u_r}(\lambda)}r
 \le\ip{v_r}{z}.
 \label{k:endpoint-support}
\end{equation}
All $u_r$ are uniformly bounded near $\lambda$; their shifted points
stay in one compact subset of $\Omega$, and the underlying subgradients
of $Q$ are bounded there.  Along a sequence $r\downarrow0$, pass to
convergent subsequences of $u_r$ and these subgradients.  Continuity of
$Q$ and $T_aQ$ makes the limiting shift active at $\lambda$, and the
closed subgradient graph makes the limiting $v_r$ a member of
$G(\lambda)$.  Thus \eqref{k:endpoint-support} implies
\[
 \ip{v_0}{z}\le\max_{v\in G(\lambda)}\ip{v}{z}.
\]
This holds in every direction.  Separation from the compact convex set
$\conv G(\lambda)$ proves $v_0\in\conv G(\lambda)$.
\end{proof}

All compactness arguments in \Cref{k:subdifferential} concern a fixed
positive $a$.  The bounds needed when $a$ decreases to zero are supplied
separately by \Cref{k:transform-properties}: active shifts, shifted
parameter compacts, and local Lipschitz constants are uniform.  At
$a=0$ the shift representation would have every $u\in\R$ active; no
compactness assertion for that degenerate representation is needed.

\section{A uniform quantitative range estimate}
\label{k:section-range}

We now compare the ideal transform with a physical attraction of finite
range.  The estimate is uniform over all preceding attraction ranges and
over arbitrarily small positive added masses.  This uniformity will allow
the error and the next mass to be chosen by a single explicit schedule.
The quadratic density term is the one underlying the Lebowitz--Penrose
limit \citep{LebowitzPenrose1966,GerardiEtAl1973}; our purpose here is a quantitative
estimate with constants uniform across the complete reference family.

\begin{theorem}[Quantitative approximation by a long-range attraction]
\label{k:range}
Fix an admissible $t$ and a compact set $\mathcal D\subset\Omega$.
There are constants $C_{\mathcal D,t}<\infty$ and
$K_{\mathcal D,t}<\infty$ with the following property.  Let $R_0\ge2$,
and let
\[
 \varphi_0=\psi_t-t\sum_j\alpha_jw_{R_j},
 \qquad 1\le R_j\le R_0,
 \qquad \alpha_j\ge0,
\]
where the sum may be empty, finite, or countably infinite.  If
$a>0$, $\sum_j\alpha_j+a\le A+1$, and $R/R_0\ge K_{\mathcal D,t}$,
then
\begin{equation}
 \sup_{\lambda\in\mathcal D}
 \abs{Q_{\varphi_0-ta w_R}(\lambda)-T_aQ_{\varphi_0}(\lambda)}
 \le C_{\mathcal D,t}\sqrt{\frac{R_0}{R}}.
 \label{k:range-estimate}
\end{equation}
Both constants are independent of $R_0$, the number and values of the
reference ranges, the reference masses, and the positive mass $a$.
\end{theorem}

Use boxes of an intermediate side $s$, with $R_0\le s\le R$.
Deleting the reference interaction between different boxes costs
$O(R_0/s)$ in pressure.  The suprema and infima of the new kernel on
pairs of boxes have row sums which, after multiplication by $s^3$,
approximate its unit integral within $O(s/R)$.  Density bounds will give
pressure errors of the same order.  The suprema will give a product upper
bound; the infima will give a lower bound from independent reference
ensembles in the boxes.  Balancing these two errors explains
the choice $s\asymp\sqrt{R_0R}$.  We begin with the deletion estimate,
which must include both the positive and negative parts of the reference
interaction.

\begin{lemma}[Grand pressure in a finite box]
\label{k:box-pressure}
Let $\varphi_0$ and $R_0$ be as in \Cref{k:range}, allowing
$\sum_j\alpha_j\le A+1$.  Write
\[
 Q^s_{\varphi_0}(\lambda)
 =\frac1{ts^3}\log\Xi_s^{\varphi_0}(\lambda).
\]
For every compact $\mathcal K\subset\Omega$ and every $s\ge R_0$ with
$s/b\in\N$,
\begin{equation}
 \sup_{\lambda\in\mathcal K}
 \abs{Q^s_{\varphi_0}(\lambda)-Q_{\varphi_0}(\lambda)}
 \le C_{\mathcal K,t}\frac{R_0}{s}.
 \label{k:box-pressure-bound}
\end{equation}
More generally, for the same $s$, partition $\Lambda_{ms}$ into $m^3$
boxes of side $s$, aligned with the small-cell grid.  In the potential
$\varphi_0-ta w_R$, where $a\ge0$, $R\ge1$, and
$\sum_j\alpha_j+a\le A+1$, delete only reference interactions between
distinct boxes.  If $\Xi^{\mathrm{del}}_{ms}$ is the resulting grand
partition function, then
\begin{equation}
 \sup_{\lambda\in\mathcal K}
 \frac{\abs{\log\Xi_{ms}^{\varphi_0-ta w_R}(\lambda)
                 -\log\Xi^{\mathrm{del}}_{ms}(\lambda)}}{t(ms)^3}
 \le C_{\mathcal K,t}\frac{R_0}{s}.
 \label{k:deletion}
\end{equation}
The constants have the uniformity asserted in \Cref{k:range} and are
also independent of $m,s,R$.
\end{lemma}

\begin{proof}
Put $S_0=t\sum_j\alpha_jw_{R_j}$ and $S_{\mathrm{new}}=ta w_R$.
The series for $S_0$ converges uniformly, is supported in the ball of
radius $R_0$, and has the row bound \eqref{f:row-bound}, including when
there are countably many reference terms.
Multiply the reference interaction on pairs belonging to distinct large
boxes by a parameter $\theta\in[0,1]$.  Every intermediate interaction
retains all intra-small-cell core terms, and its attractive part is
bounded by $S_0+S_{\mathrm{new}}$.  This has small-cell row sum at most
$M$.  Therefore \Cref{f:moments,f:perturbation} apply throughout this
interpolation, whether or not the intermediate interaction is translation
invariant.

The absolute difference is bounded on every pair by
\[
 (\psi_t+S_0)(x-y)
 \boldsymbol 1\{x,y\text{ belong to distinct large boxes}\}.
\]
Its row sum over small cells is at most
$C_b t^2+c_b t(A+1)$: the core is bounded by $t^2$ and supported within
distance two, and the attraction has the common row bound already proved.
This estimate uses an absolute majorant and thus remains valid when
$\varphi_0$ changes sign on a pair of cells.

A nonzero difference row must belong to a small cell within distance
$R_0+\sqrt3 b$ of the boundary of its large box.  Since $R_0\ge2$ and
$s\ge R_0$, there are at most $C_b s^2R_0$ such small cells per large
box.  The bound follows either by subtracting the volume of the inner
cube when the strip is thin, or from $s^3\le C s^2R_0$ when the strip
fills a fixed fraction of the box.  Thus the sum of all difference
majorants is at most
\[
 C_b m^3s^2R_0\bigl(C_b t^2+c_b t(A+1)\bigr).
\]
The finite-volume interpolation bound of \Cref{f:perturbation}, divided
by $t(ms)^3$, proves \eqref{k:deletion}.

For $a=0$ the deleted system consists of independent reference boxes,
and its grand partition function is exactly
\begin{equation}
 \Xi^{\mathrm{del}}_{ms}
       =\bigl(\Xi_s^{\varphi_0}\bigr)^{m^3}.
 \label{k:factorization}
\end{equation}
For completeness, with prescribed box counts $(n_z)$ of total $N$, the
$N!/\prod_z n_z!$ assignments of labels to boxes cancel the global
$1/N!$ and leave $\prod_z1/n_z!$.  Summing the nonnegative product
integrals over all $(n_z)$ gives \eqref{k:factorization}, including
the zero-count terms.  Its normalized logarithm equals
$Q^s_{\varphi_0}$ for every $m$.  Letting $m\to\infty$ in
\eqref{k:deletion} and using only the pressure limit of the full
reference proves \eqref{k:box-pressure-bound}.
\end{proof}

The next estimate uses the global $C^1$ bound of $w$.  Radial monotonicity
and strict positivity on its support are not required.

\begin{lemma}[Upper and lower kernel row sums]
\label{k:kernel-rows}
Let $B_z$ be the box of side $s$ centered at $z$, where $z$ ranges
over a translated lattice $s\mathbb Z^3$, and suppose $0<s\le R$.
Define
\[
 K^+_{zz'}=\sup_{x\in B_z,\,y\in B_{z'}}w_R(x-y),
 \qquad
 K^-_{zz'}=\inf_{x\in B_z,\,y\in B_{z'}}w_R(x-y).
\]
There is a constant $C_w$ depending only on $w$ and dimension such that,
with $\eta=C_ws/R$, every finite cubical collection of these boxes
satisfies
\begin{equation}
 s^3\sum_{z'}K^+_{zz'}\le1+\eta.
 \label{k:upper-row}
\end{equation}
If $z$ has distance greater than $R+4s$ from the boundary of the finite
collection, then
\begin{equation}
 s^3\sum_{z'\ne z}K^-_{zz'}\ge1-\eta.
 \label{k:lower-row}
\end{equation}
Both matrices are symmetric and nonnegative.
\end{lemma}

\begin{proof}
Symmetry follows from $w_R(x-y)=w_R(y-x)$.  Let
$G=\norm{\nabla w}_\infty$, so $\Lip(w_R)\le GR^{-4}$.
For a fixed $y\in B_{z'}$, each displacement $x-y'$ with
$x\in B_z$, $y'\in B_{z'}$ differs from $z-y$ by at most
$3s$.  Consequently
\[
 \abs{K^\pm_{zz'}-w_R(z-y)}\le3GsR^{-4}.
\]
On the infinite grid, at most $C(R/s)^3$ second boxes can have either
term nonzero: all others lie outside the support enlarged by a distance
$3s$.  Integrating over $y\in B_{z'}$ and summing thus gives
\[
 \abs{s^3\sum_{z'}K^\pm_{zz'}-
                  \int_{\R^3}w_R(z-y)\dd y}
 \le Cs^3(R/s)^3sR^{-4}=C\frac sR.
\]
The integral is one.  This comparison includes boxes meeting internal
zero regions of $w$ or crossing its support boundary, because the same
global Lipschitz estimate applies there.

Restricting the upper sum to a finite collection only decreases it.
For a center with the stated boundary distance, all potentially
contributing second boxes are still present in the lower sum.  Omitting
$z'=z$ loses at most
$s^3K^-_{zz}\le\norm w_\infty(s/R)^3$, which is bounded by
$\norm w_\infty s/R$.  Increasing the one constant $C_w$ proves both
claims.
\end{proof}

\begin{proof}[Proof of \Cref{k:range}]
We first bound the particle counts needed for the upper estimate.
For $(\beta,h)\in\mathcal D$, $s\ge1$, $V=s^3$, and $n\ge0$, put
\[
 B_n=\ee^{thn}Z^{\varphi_0}_{s,n}(\beta)
          \exp\!\left(\frac{\beta ta(1+\eta)}{2V}n^2\right),
 \qquad 0\le\eta\le1.
\]
The reference $\varphi_0$ belongs to \eqref{f:family}, and the
auxiliary quadratic coefficient satisfies $a(1+\eta)\le2(A+1)$.
Apply \Cref{f:grand-tail} on $\mathcal D$ with
$\zeta=a(1+\eta)$.  It supplies a common cutoff
$D_{\mathrm{cut}}\ge1$ and, after increasing a volume threshold
$V_0\ge1$, the bound
\begin{equation}
 \sum_{n>D_{\mathrm{cut}}V}B_n\le1
 \qquad(V\ge V_0).
 \label{k:grand-tail}
\end{equation}
Both $D_{\mathrm{cut}}$ and $V_0$ depend only on $\mathcal D,t$.
They are fixed before any local occupation constants: the tail lemma
uses only the canonical quadratic bound, even with this auxiliary
weight.

Let $U$ be the common active-shift bound on $\mathcal D$ from
\Cref{k:transform-properties}.  Choose one compact rectangle
$\mathcal D'\subset\Omega$ containing $\mathcal D$ and every point
\begin{equation}
 (\beta,h+\beta av),\qquad
 (\beta,h)\in\mathcal D,\quad
 0\le a\le A+1,\quad 0\le v\le\max\{U,D_{\mathrm{cut}}\}.
 \label{k:shifted-compact}
\end{equation}
Only now fix the moment and deletion constants of
\Cref{f:moments,k:box-pressure} on $\mathcal D'$.  All these constants
are independent of the eventual box side and of the reference.

Put $q=R/R_0$ and choose
\begin{equation}
 s=b\left\lfloor\frac{\sqrt{R_0R}}b\right\rfloor.
 \label{k:box-side}
\end{equation}
A single lower bound on $q$, uniform for $R_0\ge2$, ensures
\begin{equation}
 R_0\le s\le R,\qquad
 \frac12R_0\sqrt q\le s\le R_0\sqrt q,\qquad
 s^3\ge V_0,\qquad
 \eta=C_w\frac sR\le1.
 \label{k:ratio-conditions}
\end{equation}
Partition $\Lambda_L$, $L=ms$, into $m^3$ boxes $B_z$, and delete
reference interactions between different boxes as in
\Cref{k:box-pressure}.  The new attraction remains intact.  All limits
in $m$ below hold with $s,R,R_0$, the reference, and the parameters fixed.

\medskip\noindent\emph{Upper bound.}
Let $N_z$ be the particle count in $B_z$, and write
$W_R=\sum_{i<j}w_R(x_i-x_j)$.  Adding nonnegative diagonal self terms,
then using symmetry and $2N_zN_{z'}\le N_z^2+N_{z'}^2$, gives
\begin{align*}
 W_R
 &\le\frac12\sum_{z,z'}K^+_{zz'}N_zN_{z'}
 \le\frac12\sum_zN_z^2\sum_{z'}K^+_{zz'}
 \le\frac{1+\eta}{2s^3}\sum_zN_z^2.
\end{align*}
Splitting labels among boxes as in \eqref{k:factorization}, we obtain
\begin{equation}
 \Xi_L^{\mathrm{del}}(\beta,h)
 \le\left[\sum_{n\ge0}
 \ee^{thn}Z^{\varphi_0}_{s,n}(\beta)
 \exp\!\left(\frac{\beta ta(1+\eta)}{2s^3}n^2\right)
 \right]^{m^3}.
 \label{k:upper-product}
\end{equation}
For $0\le n\le D_{\mathrm{cut}}s^3$, let $d=n/s^3$.  The canonical term at count
$n$ is one of the nonnegative terms in the reference grand sum at the
shifted field $h+\beta ad$.  Therefore
\begin{align*}
 \frac1{ts^3}\log B_n
 &\le Q^s_{\varphi_0}(\beta,h+\beta ad)
                     -\frac{\beta a}{2}d^2
                     +\frac{\beta a\eta}{2}d^2\\
 &\le Q_{\varphi_0}(\beta,h+\beta ad)
                     -\frac{\beta a}{2}d^2
                     +C_{\mathcal D,t}\left(\frac{R_0}{s}+\eta\right)\\
 &\le T_aQ_{\varphi_0}(\beta,h)
                     +C_{\mathcal D,t}\left(\frac{R_0}{s}+\eta\right).
\end{align*}
All shifted parameters lie in $\mathcal D'$ by
\eqref{k:shifted-compact}, so the finite-box comparison applies; the bound
$d\le D_{\mathrm{cut}}$ controls the extra quadratic.  The argument includes $n=0$,
whose term equals one.  There are $\lfloor D_{\mathrm{cut}}s^3\rfloor+1$ such
counts; the remaining terms have total at most one by
\eqref{k:grand-tail}.  Since $T_aQ_{\varphi_0}\ge0$, taking logarithms
in \eqref{k:upper-product} yields
\begin{align}
 \frac1{tL^3}\log\Xi_L^{\mathrm{del}}
 \le T_aQ_{\varphi_0}
 +C_{\mathcal D,t}\left(\frac{R_0}{s}+\eta\right)
 +\frac{\log(\lfloor D_{\mathrm{cut}}s^3\rfloor+2)}{ts^3}.
 \label{k:upper-final}
\end{align}
Combine this with \eqref{k:deletion} and let $m\to\infty$ in the
full physical pressure.  This proves the desired upper estimate with
the three displayed errors.

\medskip\noindent\emph{Lower bound.}
Fix an active shift $u\in[0,U]$ for $T_aQ_{\varphi_0}(\beta,h)$.
Let $\Pi$ be the product of the $m^3$ reference grand measures in the
boxes, all at $(\beta,h+\beta au)$, and put
\[
 d_s=s^{-3}\mathbb E_\Pi N_z.
\]
The common small-cell moment bound on $\mathcal D'$ gives a uniform
bound $d_s\le C_{\mathcal D,t}$.  Indeed each box consists of exactly
$(s/b)^3$ small cells, and Cauchy--Schwarz bounds each mean count by the
square root of its uniformly bounded second moment.

The change of field and insertion of the new attraction give the exact
identity
\begin{equation}
 \Xi_L^{\mathrm{del}}(\beta,h)
 =\bigl[\Xi_s^{\varphi_0}(\beta,h+\beta au)\bigr]^{m^3}
       \mathbb E_\Pi
       \exp\{\beta ta(W_R-uN)\},
 \label{k:jensen-identity}
\end{equation}
where $N=\sum_zN_z$.  The box factorials in $\Pi$ are precisely those
obtained by splitting the full grand measure among boxes.  All terms
are finite by the common quadratic stability bound.  Jensen's inequality
therefore gives
\begin{equation}
 \frac1{tL^3}\log\Xi_L^{\mathrm{del}}
 \ge Q^s_{\varphi_0}(\beta,h+\beta au)
       -\beta au d_s+\frac{\beta a}{L^3}\mathbb E_\Pi W_R.
 \label{k:jensen-lower}
\end{equation}

Discard the nonnegative contributions to $W_R$ from pairs in one box.
For the remaining pairs the kernel infima give
\[
 W_R\ge\frac12\sum_{z\ne z'}K^-_{zz'}N_zN_{z'}.
\]
Entire distinct boxes are independent under $\Pi$, so for $z\ne z'$,
$\mathbb E_\Pi[N_zN_{z'}]=s^6d_s^2$.  No independence inside a box is
used.  Let $m_{\mathrm{int}}$ be the number of centers with distance
greater than $R+4s$ from the outer boundary.  Nonnegativity of the other
rows and \Cref{k:kernel-rows} imply
\[
 \frac1{L^3}\mathbb E_\Pi W_R
 \ge\frac{m_{\mathrm{int}}}{m^3}\,
                    \frac{1-\eta}{2}d_s^2.
\]
For fixed $s,R$, $m_{\mathrm{int}}/m^3\to1$ as $m\to\infty$.
Consequently \eqref{k:jensen-lower} and \eqref{k:deletion}, using the
known limit of the full physical model, show that
\begin{align*}
 Q_{\varphi_0-ta w_R}(\beta,h)
 &\ge Q^s_{\varphi_0}(\beta,h+\beta au)
       +\beta a\left(-u d_s+\frac{1-\eta}{2}d_s^2\right)
       -C_{\mathcal D,t}\frac{R_0}{s}.
\end{align*}
The exact square identity
\begin{equation}
 -u d_s+\frac{1-\eta}{2}d_s^2
 =-\frac{u^2}{2}+\frac{(d_s-u)^2}{2}-\frac\eta2d_s^2
 \ge-\frac{u^2}{2}-C_{\mathcal D,t}\eta
 \label{k:mean-square}
\end{equation}
and \eqref{k:box-pressure-bound} at the shifted parameter now give
\begin{align}
 Q_{\varphi_0-ta w_R}(\beta,h)
 &\ge Q_{\varphi_0}(\beta,h+\beta au)-\frac{\beta a}{2}u^2
       -C_{\mathcal D,t}\left(\frac{R_0}{s}+\eta\right)\notag\\
 &=T_aQ_{\varphi_0}(\beta,h)
       -C_{\mathcal D,t}\left(\frac{R_0}{s}+\eta\right).
 \label{k:lower-final}
\end{align}
In particular no convergence of $d_s$ to $u$, or convergence of $d_s$
at all, is required.  This proves the same bound at reference coexistence
parameters.

Finally, \eqref{k:ratio-conditions} gives
$R_0/s\le2q^{-1/2}$ and $s/R\le q^{-1/2}$.  For $s\ge1$,
\[
 \frac{\log(\lfloor D_{\mathrm{cut}}s^3\rfloor+2)}{ts^3}
 \le\frac{C_{\mathcal D,t}}s
 \le C_{\mathcal D,t}q^{-1/2},
\]
where the last step uses $R_0\ge2$ and
$s\ge\tfrac12R_0\sqrt q$.  Equations
\eqref{k:upper-final} and \eqref{k:lower-final}, with the deletion cost,
prove \eqref{k:range-estimate}.

No step imposes a positive lower bound on $a$, and no error has been
divided by $a$.  Only the reference and full physical models require
thermodynamic limits; the deleted systems were treated by exact
factorization or finite-volume inequalities.
\end{proof}

\section{An initial triple of pressure slopes}
\label{i:section}

The narrow zero-energy shell in the core permits clusters of four particles,
with relative volume of order \(l_t^9=\ee^{-9t}\).  This entropy cost,
together with the packing density \(p\), produces three competing branches
after the first ideal attraction.  We show that, for one sufficiently large
finite \(t\), the pressure subdifferential throughout a parameter neighborhood
is generated by actual supports near the three limiting slopes.  The
subsequent iteration will preserve this separation and find a parameter
where all three groups are present.

In this section alone we let \(t\) tend to infinity to obtain the required
estimates; we fix its finite value at the end.

\subsection{Clusters and the bare-core pressure}

\begin{proposition}[Pressure of the bare core]
\label{i:bare-pressure}
Let \(q^{(t)}=Q_{\psi_t}\), with the pressure normalization and the core
of \Cref{f:thermodynamics,f:core}.  Then, locally uniformly on
\(\Omega=(0,\infty)\times\R\),
\begin{equation}
 q^{(t)}(\beta,h)\longrightarrow
 q_*(\beta,h):=\max\{0,ph,4ph-9p\}.
 \label{i:bare-limit}
\end{equation}
\end{proposition}

\begin{proof}
Fix a compact parameter set \(K\subset\Omega\), and choose
\(\beta_->0\) and \(H<\infty\) such that
\(\beta\ge\beta_-\) and \(\abs{h}\le H\) on \(K\).
Throughout the proof \(V=L^3\), and
\(I_t=[l_t,(1+\Delta)l_t]\).

\emph{Free configurations and the upper bound.}
Call a pair bad when its distance is less than one and is not in \(I_t\).
A configuration is free if it has no bad pair.  In a free configuration,
join distinct particles whose distance is less than one.  Every edge then
has length at most \((1+\Delta)l_t\).  The endpoints of a path of two
edges are therefore at distance at most \(2(1+\Delta)l_t<1\), and hence
are joined.  Each connected component is consequently a clique with all
pair distances in \(I_t\).  By \Cref{f:shell-graph}, its size is at most
four.  Choosing one representative
from every component gives a set with pairwise distances at least one,
so the component count is at most \(M_L\).

Consider the grand integral over free configurations with their energy
discarded:
\[
 \Xi_{L,\mathrm{free}}(h)
 =\sum_{n\ge0}\frac{\ee^{thn}}{n!}
       \int_{\Lambda_L^n}\mathbf 1_{\mathrm{free}}(x_1,\ldots,x_n)
          \dd x_1\cdots\dd x_n.
\]
For an ordered list of component sizes \(m_1,\ldots,m_k\), with
\(m_i\in\{1,2,3,4\}\) and \(\sum_i m_i=n\), the number of ordered
partitions of the particle labels into these blocks is
\(n!/\prod_i m_i!\).  On summing over all ordered size lists, division
by \(k!\) counts each unordered family of label blocks once.  This
remains true when sizes repeat: the label blocks themselves are distinct
and have exactly \(k!\) orderings.

In a block of size \(m\), choose its least label as the root.  Its position has volume
at most \(V\), and each of the other \(m-1\) particles lies in a ball of
volume \(C l_t^3\) about it, where \(C=4\pi(1+\Delta)^3/3\).
After cancellation of the global factor \(1/n!\), the block therefore
contributes at most \(V(C l_t^3)^{m-1}/m!\).  Discarding all further
restrictions within and between components and summing their sizes gives
\begin{equation}
 \Xi_{L,\mathrm{free}}(h)
 \le\sum_{k=0}^{M_L}\frac{(V\Theta_t(h))^k}{k!},
 \qquad
 \Theta_t(h)=\sum_{m=1}^4
       \frac{\ee^{thm}(C l_t^3)^{m-1}}{m!}.
 \label{i:rooted-count}
\end{equation}

To pass from free configurations to all configurations, let \(B(X)\)
be the set of labels incident to a bad pair.  There are at least
\(\abs{B(X)}/2\) bad pairs, and each has core energy exactly \(t^2\).
All other core terms are nonnegative, and \(X\setminus B(X)\) is free.
Thus the following pointwise bound holds:
\begin{equation}
 \ee^{-\beta U^{\psi_t}(X)}
 \le\sum_{J\subset\{1,\ldots,n\}}
       \ee^{-\beta t^2\abs{J}/2}
       \mathbf 1_{\mathrm{free}}(X\setminus J).
 \label{i:removal-pointwise}
\end{equation}
Indeed, the summand \(J=B(X)\) already bounds the left side.
For \(r=\abs{J}\) and \(m=n-r\), the label factor is
\(\binom{m+r}{r}/(m+r)!=1/(r!m!)\).  The removed positions integrate
to \(V^r\).  Tonelli's theorem and the grand sum now give the exact
factorization of the upper bound
\begin{equation}
 \Xi_L^{\psi_t}(\beta,h)
 \le
 \exp\!\left(V\ee^{th-\beta t^2/2}\right)
       \Xi_{L,\mathrm{free}}(h).
 \label{i:removal-factor}
\end{equation}
Shell endpoints are included in \(I_t\), although their energy may be
positive.  Pairs at distance exactly one are not edges of the component
graph.  Both conventions merely discard nonnegative energy in the upper
bound, so these arguments hold pointwise, including at the endpoints.

Set \(a_t(h)=t^{-1}\log\Theta_t(h)\).  The four exponential rates in
\(\Theta_t\) are
\[
 h,\qquad 2h-3,\qquad 3h-6,\qquad 4h-9.
\]
The middle two are convex combinations of the first and last.  Their
positive prefactors in \eqref{i:rooted-count} are fixed, so
\begin{equation}
 a_t(h)=\max\{h,4h-9\}+O(t^{-1}),
 \label{i:cluster-rates}
\end{equation}
uniformly for \(\abs{h}\le H\).  For \(k\ge1\), the inequality
\(k!\ge(k/\ee)^k\), with \(x=k/V\), gives
\[
 \frac1{tV}\log\frac{(V\Theta_t)^k}{k!}
 \le x a_t(h)+\frac{x(1-\log x)}t
 \le \frac{M_L}{V}\max\{0,a_t(h)\}+\frac1t,
\]
since \(x(1-\log x)\le1\) for \(x>0\).  The same final bound covers
\(k=0\).  There are only \(M_L+1=O(1+V)\) terms in
\eqref{i:rooted-count}.  First letting \(L\to\infty\), with \(t\)
fixed, and using \Cref{f:packing,f:thermodynamics}, we obtain
\begin{equation}
 q^{(t)}(\beta,h)
 \le p\max\{0,a_t(h)\}+\frac1t
       +\frac1t\ee^{th-\beta t^2/2}.
 \label{i:upper-pressure}
\end{equation}
The final term is at most
\(t^{-1}\ee^{tH-\beta_-t^2/2}\) on \(K\).  Thus
\eqref{i:cluster-rates} proves the required uniform upper bound.

\emph{Zero-energy configurations and the lower bound.}
Fix \(s>0\) and \(0<\nu\le1\).  Dilate a maximizing packing in a
side-\(s\) cube by \(1+\nu\), and center it in a cube of side
\[
 P=(s+2)(1+\nu).
\]
Repeat this pattern in adjacent period cubes.  Centers in one period
have separation at least \(1+\nu\); centers in different periods have
at least this separation as well.  Every center has distance at least
\(1+\nu\) from the boundary of its period cube.  A cube of side \(mP\)
contains exactly \(k=m^3M_s\) centers.

Choose \(r\in(0,\nu/16)\) and let \(b_0=4\pi r^3/3\).
At every center \(z\), put the root particle in \(B(z,r)\).
There are two constructions, with cluster size \(j=1\) or \(j=4\).
For \(j=4\), choose vectors \(a_2,a_3,a_4\in\R^3\) such that
\(0,a_2,a_3,a_4\) are the vertices of a regular tetrahedron of side
\(1+\Delta/2\).  Fix \(\delta\in(0,\Delta/24)\) and require
\[
 x_i-x_1\in B(l_t a_i,\delta l_t),\qquad i=2,3,4.
\]
A root-to-vertex distance differs from
\((1+\Delta/2)l_t\) by at most \(\delta l_t\), and another edge
differs by at most \(2\delta l_t\).  All six distances consequently
lie strictly inside
\[
 ((1+\Delta/3)l_t,(1+2\Delta/3)l_t),
\]
where the core vanishes.  Write \(b_1=4\pi\delta^3/3>0\).
The change from positions to the root and the three relative vectors
has Jacobian one.  For one prescribed label ordering the integration
volume at each center is therefore
\(b_0(b_1l_t^3)^3\), with positive relative volume \(b_1^3l_t^9\).
For \(j=1\), the corresponding volume is \(b_0\).

Every particle lies within \(r+2l_t\) of its center.  For all sufficiently
large \(t\), depending on \(\nu,r\) but not on \(m,\beta,h\), these
center neighborhoods are disjoint and lie inside their period cubes,
and intercluster distances are greater than
\(1+\nu/2>1+1/t\).  Hence all configurations just described have
exactly zero core energy.

The centers are distinguishable.  Assign an unordered set of \(j\)
particle labels to each of the \(k\) centers, and retain just one
deterministic ordering of each assigned set, for example the increasing
order.  There are \((jk)!/(j!)^k\) such assignments.  Different
assignments send some label to different disjoint center neighborhoods,
so their integration events are disjoint.  Possible alternative choices
of the root within the same set are never summed.  Dividing by the
canonical factorial \((jk)!\) gives
\begin{equation}
 Z^{\psi_t}_{mP,jk}(\beta)
 \ge\left(\frac{b_0(b_1l_t^3)^{j-1}}{j!}\right)^k,
 \qquad j\in\{1,4\}.
 \label{i:cluster-lower}
\end{equation}
Using this term in the grand sum and first letting \(m\to\infty\), we
find
\begin{equation}
 q^{(t)}(\beta,h)
 \ge c_{s,\nu}\left(jh-3(j-1)
       +\frac1t\log\frac{b_0b_1^{j-1}}{j!}\right),
 \qquad
 c_{s,\nu}=\frac{M_s}{((s+2)(1+\nu))^3}.
 \label{i:lower-pressure}
\end{equation}
The vacuum also gives \(q^{(t)}\ge0\).  The constants in
\eqref{i:lower-pressure} have no dependence on \(\beta,h\), and
\(c_{s,\nu}\) can be made arbitrarily close to \(p\) by taking \(s\)
large and \(\nu\) small.  More explicitly, on \(\abs{h}\le H\) the
two functions \(jh-3(j-1)\), \(j=1,4\), have a common finite absolute
bound.  Given any error tolerance, first choose \(s,\nu\) so that
replacing \(c_{s,\nu}\) by \(p\) changes both branches by less than
that tolerance, and then take \(t\) large enough to make the two
logarithmic terms small.  Taking the maximum of these branches and zero
preserves the uniform error, even where a nonvacuum branch is negative.
Together with \eqref{i:upper-pressure}, this proves
\eqref{i:bare-limit} on \(K\).
\end{proof}

\subsection{The limiting attraction and its active densities}

We next apply the ideal attraction with mass \(A=3/(2p)\).  Let
\(P_1^t=T_Aq^{(t)}\), and define
\begin{equation}
 \begin{split}
 L_1(\beta,h)&=0,\\
 L_2(\beta,h)&=ph+\frac{\beta A p^2}{2},\\
 L_3(\beta,h)&=4ph-9p+\frac{\beta A(4p)^2}{2},\\
 P_*(\beta,h)&=\max_{1\le i\le3}L_i(\beta,h),
 \qquad D_*=\{0,p,4p\}.
 \end{split}
 \label{i:three-planes}
\end{equation}

\begin{lemma}[The limiting transform and its optimizer gap]
\label{i:limiting-transform}
The functions \(P_1^t\) converge locally uniformly on \(\Omega\) to
\(P_*\).  For every compact \(K\subset\Omega\), all their active
optimizers lie in one interval \([0,U]\) for all sufficiently large
\(t\), and
\[
 \sup_{\lambda\in K}\ \sup_{u\ \mathrm{active\ for}\ P_1^t(\lambda)}
       \dist(u,D_*)\longrightarrow0.
\]
In particular, this holds at parameters where several limiting planes
meet.  If \(\beta_-=\min_{(\beta,h)\in K}\beta\), the limiting
objective obeys the explicit bound
\begin{equation}
 F_*(\lambda,u):=q_*(\beta,h+\beta Au)-\frac{\beta A u^2}{2}
 \le P_*(\lambda)-\frac{\beta_-A}{2}\dist(u,D_*)^2
 \quad(\lambda\in K,\ u\in\R).
 \label{i:objective-gap}
\end{equation}
\end{lemma}

\begin{proof}
Apply the uniform-in-\(t\) bound in \Cref{k:transform-properties} to
the bare core \(\psi_t\), mass \(A\), and parameter compact \(K\).
For all sufficiently large \(t\), every active shift lies in one fixed
interval \([0,U]\).  Enlarge \(U\), if necessary, so that \(4p\le U\).

Write \(d_1=0,d_2=p,d_3=4p\).  Completing the square gives the identity
\begin{equation}
 F_*(\lambda,u)
 =\max_{1\le i\le3}
       \left\{L_i(\lambda)-\frac{\beta A}{2}(u-d_i)^2\right\}.
 \label{i:limiting-parabolas}
\end{equation}
Its supremum over \(u\) is \(P_*(\lambda)\), and all its maximizing
values belong to \(D_*\).  The same identity proves
\eqref{i:objective-gap}.

The shifted points
\[
 K'=\{(\beta,h+\beta Au):(\beta,h)\in K,\ 0\le u\le U\}
\]
form a compact subset of \(\Omega\).  Put
\(E_t=\sup_{K'}\abs{q^{(t)}-q_*}\), which tends to zero by
\Cref{i:bare-pressure}.  Both transforms can be maximized over
\([0,U]\), so
\[
 \sup_K\abs{P_1^t-P_*}\le E_t.
\]
If \(u\) is active for \(P_1^t(\lambda)\), approximation of its
objective and of the maximum shows
\(F_*(\lambda,u)\ge P_*(\lambda)-2E_t\).  By
\eqref{i:objective-gap},
\begin{equation}
 \dist(u,D_*)
 \le 2\sqrt{\frac{E_t}{\beta_-A}}.
 \label{i:active-density-error}
\end{equation}
This estimate is uniform on \(K\) and requires no strict inequality
between the three affine branches.
\end{proof}

\subsection{Uniform separation of generating slopes}

The optimizer gap locates the density coordinates of active supports.
We now identify neighborhoods of the full slopes that will remain separated
under the iteration.  The slopes of the three limiting planes and their
barycenter are
\begin{equation}
 \begin{gathered}
 d_1=0,\qquad d_2=p,\qquad d_3=4p,\qquad
 v_i=(Ad_i^2/2,d_i),\\
 v_1=(0,0),\qquad v_2=(3p/4,p),\qquad v_3=(12p,4p),\\
 g=\frac{v_1+v_2+v_3}{3}=(17p/4,5p/3),\qquad
 \lambda_*=(1,-3/4).
 \end{gathered}
 \label{i:slopes}
\end{equation}
All three planes vanish at \(\lambda_*\), and
\(\det(v_2-v_1,v_3-v_1)=-9p^2\ne0\).  Thus \(g\) is strictly
inside their slope triangle.

\begin{lemma}[Robust slope geometry]
\label{i:geometry}
There are constants \(0<\tau<p/16\) and \(\gamma>0\) with the
following properties.  Every choice
\(w_i\in\clB{v_i}{\tau}\), \(i=1,2,3\), satisfies
\begin{equation}
 \clB{g}{\gamma}\subset\conv\{w_1,w_2,w_3\}.
 \label{i:robust-triangle}
\end{equation}
For each distinct pair \(i,j\),
\begin{equation}
 g\notin\conv\bigl(\clB{v_i}{\tau}\cup\clB{v_j}{\tau}\bigr).
 \label{i:two-groups}
\end{equation}
Moreover, for every \(\lambda\in\Omega\),
\begin{equation}
 P_*(\lambda)-P_*(\lambda_*)
       -\ip{g}{\lambda-\lambda_*}
 \ge\gamma\norm{\lambda-\lambda_*}.
 \label{i:linear-growth}
\end{equation}
\end{lemma}

\begin{proof}
Because \(g\) is an interior point of the nondegenerate triangle, the
two numbers
\[
 \eta=\min_{\norm{z}=1}\max_i\ip{v_i-g}{z}>0,
 \qquad
 \zeta=\min_{i<j}\dist\bigl(g,[v_i,v_j]\bigr)>0
\]
are positive.  For the first assertion, one may also see positivity
directly: if all three inner products were nonpositive, their sum would
be zero, so all would vanish, contradicting noncollinearity.
Choose
\[
 0<\tau<\min\{p/16,\eta/4,\zeta/2\},
 \qquad \gamma=\eta/2.
\]
For any unit vector \(z\) and any allowed vertices \(w_i\),
\[
 \max_i\ip{w_i-g}{z}\ge\eta-\tau>\gamma.
\]
The support-function characterization of inclusion of compact convex
sets proves \eqref{i:robust-triangle}.  Every point of the hull of two
radius-\(\tau\) balls lies within distance \(\tau\) of the segment
joining their centers.  Since \(\tau<\zeta\), this also proves
\eqref{i:two-groups}.  Finally, all three planes agree at
\(\lambda_*\), and hence for \(z=\lambda-\lambda_*\),
\[
 P_*(\lambda)-P_*(\lambda_*)-\ip{g}{z}
 =\max_i\ip{v_i-g}{z}\ge\eta\norm{z}\ge\gamma\norm{z}.
\]
\end{proof}

The geometry of \Cref{i:geometry} is illustrated in
\Cref{fig:support-geometry}.
\begin{figure}[tb]
\centering
\begin{tikzpicture}[>=Latex,font=\small]
  \begin{scope}[scale=.48]
    \draw[->,black!55] (-.45,0) -- (13.1,0)
      node[below=2pt,text=black] {\(e/p\)};
    \draw[->,black!55] (0,-.35) -- (0,5.0)
      node[above,text=black] {\(d/p\)};
    \fill[blue!8] (0,0) -- (.75,1) -- (12,4) -- cycle;
    \draw[blue!60!black,thick] (0,0) -- (.75,1) -- (12,4) -- cycle;
    \fill (0,0) circle (2.4pt) node[below left] {\(v_1\)};
    \fill (.75,1) circle (2.4pt) node[above left] {\(v_2\)};
    \fill (12,4) circle (2.4pt) node[above right] {\(v_3\)};
    \fill[orange!75!black] (4.25,1.666667) circle (2.4pt);
    \draw[orange!85!black,thick] (4.25,1.666667) circle (.115);
    \node[above left] at (4.0,2.25) {\(g\)};
    \draw[black!55] (3.93,2.12) -- (4.18,1.80);
    \draw[black!55,dashed] (3.75,1.15) rectangle (4.75,2.2);
  \end{scope}
  \draw[black!55,->] (2.28,1.05) to[bend left=14] (7.48,2.3);
  \begin{scope}[shift={(8.9,1.65)}]
    \fill[orange!8] (0,0) circle (1.35);
    \draw[orange!85!black,thick] (0,0) circle (1.35);
    \draw[black!55,dashed] (-1.7,0) -- (1.7,0);
    \draw[very thick,blue!65!black] (-.93,0) -- (.93,0);
    \fill[blue!65!black] (-.93,0) circle (2pt)
      node[below=5pt] {\((s_1,\rho)\)};
    \fill[blue!65!black] (.93,0) circle (2pt)
      node[below=5pt] {\((s_2,\rho)\)};
    \fill (0,0) circle (1.6pt) node[above=4pt] {\(g\)};
    \node at (0,1.67) {\(\overline B(g,\gamma)\)};
    \node at (0,-1.72) {\(\rho=5p/3,\quad s_1<s_2\)};
  \end{scope}
\end{tikzpicture}
\caption{Pressure-support coordinates \((e,d)\), with the left panel
drawn in coordinates \((e/p,d/p)\).
The triangle of the initial limiting slopes \(v_i\) contains a ball
about its barycenter \(g\);
the ball persists under the controlled perturbations of its vertices.
The schematic inset shows two points on the central horizontal slice
at density \(5p/3\).  Once the ball lies in the subdifferential of a
pressure at one parameter, each nearby horizontal slice supplies two
supports with distinct inverse-temperature slopes at that same parameter.
Lemma~\ref{m:canonical-corner} converts this geometry into the canonical
derivative jump throughout a density interval.
}
\label{fig:support-geometry}
\end{figure}


Define the closed parameter neighborhoods
\begin{equation}
 W(r)=\clB{\lambda_*}{r},\qquad W_1=W(1/4)\subset\Omega.
 \label{i:neighborhoods}
\end{equation}
The next proposition gives a statement about actual supporting slopes
throughout \(W_1\), as required for the subsequent iteration.

\begin{proposition}[Initial separation of generating slopes]
\label{i:initial-separation}
For all sufficiently large finite \(t\), the size conditions in
\Cref{f:stability,k:transform-properties} hold and
\begin{equation}
 \sup_{\lambda\in W_1}\abs{P_1^t(\lambda)-P_*(\lambda)}
       <\frac{\gamma}{64}.
 \label{i:initial-pressure-error}
\end{equation}
At every \(\lambda\in W_1\) one also has the exact identity
\begin{equation}
 \partial P_1^t(\lambda)
 =\conv\left(
       \partial P_1^t(\lambda)\cap
       \bigcup_{i=1}^3\clB{v_i}{\tau/2}\right).
 \label{i:initial-hull}
\end{equation}
Fix one such \(t\) for the rest of the construction.
\end{proposition}

\begin{proof}
Use the uniform optimizer bound \([0,U]\) from
\Cref{i:limiting-transform} with \(K=W_1\).  All shifted points
\(x=(\beta,h+\beta Au)\) lie in the fixed compact \(K'\) from that
proof.  Choose one \(\sigma>0\) such that
\[
 K''=\{x+(s,0):x\in K',\ \abs{s}\le\sigma\}\subset\Omega,
\]
and set \(E'_t=\sup_{K''}\abs{q^{(t)}-q_*}\to0\).
For every \((e,d)\in\partial q^{(t)}(x)\), its two supporting
inequalities at \(x+(\sigma,0)\) and \(x-(\sigma,0)\), together
with the independence of \(q_*\) from \(\beta\), imply
\begin{equation}
 -\frac{2E'_t}{\sigma}\le e\le\frac{2E'_t}{\sigma}.
 \label{i:energy-support-error}
\end{equation}
These are fixed steps on an enlarged compact, so the estimate also
covers outer parameters on the boundary of \(W_1\).

At a shifted point corresponding to an active \(u\),
\Cref{k:active-parabola} gives \(d=u\) for every underlying support.
By \Cref{k:subdifferential}, the actual active generators of
\(\partial P_1^t(\lambda)\) are precisely the vectors in the union
\[
 G_t(\lambda)=
 \bigcup_{u\ \mathrm{active}}
   \left\{(e+Au^2/2,u):(e,u)\in
           \partial q^{(t)}(\beta,h+\beta Au)\right\},
 \qquad
 \partial P_1^t(\lambda)=\conv G_t(\lambda).
\]
For a nearest \(d_i\in D_*\), every such generator \(w\)
satisfies
\[
 \norm{w-v_i}
 \le\frac{2E'_t}{\sigma}
       +\left(1+\frac A2(U+4p)\right)\abs{u-d_i}.
\]
Equations \eqref{i:active-density-error} and
\eqref{i:energy-support-error} show that this tends uniformly to zero
over all outer parameters, all active optimizers, and all underlying
supports.  Thus, for all sufficiently large \(t\),
\[
 G_t(\lambda)\subset
       \partial P_1^t(\lambda)\cap
       \bigcup_i\clB{v_i}{\tau/2}
       \qquad(\lambda\in W_1).
\]
Taking convex hulls proves one inclusion in \eqref{i:initial-hull};
the other follows from convexity of the subdifferential.  Finally,
\eqref{i:initial-pressure-error} follows from
\Cref{i:limiting-transform}.  All earlier size conditions are lower
bounds on \(t\), so they can be imposed at the same time.
\end{proof}

\section{The multiscale potential and the canonical singularity}
\label{sec:iteration}

Fix \(t\) as in \Cref{i:initial-separation}.  We will preserve its
three groups of generating supports while replacing ideal attractions
by physical ones.  The essential step is quantitative: a pressure error
\(\delta\), followed by an ideal attraction of mass \(a\), enlarges
the slope groups by \(O(\delta/a+a)\) and costs \(O(a)\) in the
parameter domain.  We first prove this step, then choose one sequence
of masses and ranges for which both losses are summable.

\subsection{Preservation under one attraction}

Let \(q\) be a physical pressure and \(P\) a convex approximation
whose subdifferential is generated near the three initial slopes.
In the iteration, \(P\) will be the preceding ideal transform and
\(q\) its finite-range realization.  We ask how closely the generating
supports of \(T_aq\) remain in those three groups when
\(\norm{P-q}_{W_1}\le\delta\).  At each active shifted point, the
parabolic bound first isolates one group in a neighborhood of radius
proportional to \(a\).  Secant inequalities in that neighborhood then
transfer the support bound from \(P\) to \(q\).

Recall the parameter balls
\(W(r)=\clB{\lambda_*}{r}\), \(W_1=W(1/4)\),
the slope vectors \(v_i=(Ad_i^2/2,d_i)\), and the constants
\(\tau,\gamma>0\) from \Cref{i:geometry}.
By \Cref{k:transform-properties}, choose \(U_*>0\) bounding all active
shifts for every transform \(T_aQ\) of a family pressure, with \(0<a\le A+1\), at parameters
in \(W_1\), uniformly over admissible references and masses.
Uniform convex-function bounds on a larger compact neighborhood give a
common Lipschitz constant \(C_*\ge1\) for all these transforms on \(W_1\).
These constants depend only on the fixed family and parameter compacts.
Put \(\beta_{\min}=3/4\), \(\beta_{\max}=5/4\).
Choose \(k>0\), and then \(l>0\), sufficiently small that
\begin{equation}\label{eq:band-constants}
 \frac{k}{2\beta_{\min}}+\frac{2C_*l}{k}<\frac p{16}.
\end{equation}

Set
\[
 S_* = \beta_{\max}U_*+k+l+1.
\]
Here \(\beta_{\max}U_*a\) bounds the displacement to an active
shifted point, \(ka\) will be a field increment, and \(la\) the
radius on which all supports will be controlled.  All these constants
depend only on the fixed \(t\), the family, and the parameter compacts.

\begin{lemma}[One-step preservation of generating supports]
\label{m:one-step}
Let \(q\) be a pressure from the family \eqref{f:family}, and let
\(P\) be a finite convex function on \(\Omega\) that is
\(C_*\)-Lipschitz on \(W_1\).  Suppose
\[
 \begin{gathered}
 0<a\le A+1,\qquad 0<r<\tau,\qquad S_*a<s\le\frac14,\\
 \delta\ge0,\qquad \norm{P-q}_{W_1}\le\delta,\qquad
 \frac{2\delta}{ka}<\frac p{16}.
 \end{gathered}
\]
Assume that throughout \(W(s)\),
\begin{equation}\label{eq:one-step-hypothesis}
 \partial P(y)=\conv\left(\partial P(y)\cap
                 \bigcup_{i=1}^3\clB{v_i}{r}\right).
\end{equation}
Then, with
\[
 r^+=r+\frac{2\delta}{la}+\frac12aU_*^2,
\]
one has for every \(\lambda\in W(s-S_*a)\),
\begin{equation}\label{eq:one-step-conclusion}
 \partial(T_aq)(\lambda)
 =\conv\left(\partial(T_aq)(\lambda)\cap
                 \bigcup_{i=1}^3\clB{v_i}{r^+}\right).
\end{equation}
\end{lemma}

\begin{proof}
Fix \(\lambda=(\beta,h)\in W(s-S_*a)\), and choose an active
shift \(u\) for \(T_aq\) at \(\lambda\).
Let \(x=(\beta,h+\beta au)\).
The distance from \(\lambda\) to the boundary of \(W(s)\) is at least
\(S_*a\).
Since \(\abs{x-\lambda}\le\beta_{\max}U_*a\), every point
\[
 y,\quad y+(0,\pm ka),\quad x+(0,\pm ka),
 \qquad \abs{y-x}\le la,
\]
lies in \(W(s)\).
The surplus \(a\) leaves these test points in the interior of \(W(s)\).

\emph{A common density band.}
For \((e,d)\in\partial P(y)\), the supporting inequality at \(y\),
the Lipschitz bound on \(P\), and
the bound \(\norm{P-q}_{W_1}\le\delta\) give
\begin{align*}
 \pm ka\,d
 &\le P(y+(0,\pm ka))-P(y)\\
 &\le q(x+(0,\pm ka))-q(x)
                         +2\delta+2C_*la\\
 &\le \pm ka\,u+\frac{(ka)^2}{2\beta a}
                         +2\delta+2C_*la.
\end{align*}
The last step is the active-point parabola of \Cref{k:active-parabola}.
Using both signs, \eqref{eq:band-constants}, and the hypothesis on
\(\delta/a\), we obtain
\begin{equation}\label{eq:common-density-band}
 \abs{d-u}\le
 \frac{k}{2\beta_{\min}}+\frac{2C_*l}{k}
                       +\frac{2\delta}{ka}<\frac p8.
\end{equation}
The same \(u\) works for every \(y\) in this neighborhood.

\emph{One slope group near the shifted point.}
The set in the right side of \eqref{eq:one-step-hypothesis} is nonempty:
the full subdifferential is nonempty, and the assumed hull identity
identifies it with the convex hull of that set.
If a generator belongs to \(\clB{v_i}{r}\), then
\[
 \abs{u-d_i}<\frac p8+\tau<\frac{3p}{16}.
\]
Since the distinct \(d_i\) are separated by at least \(p\), at most one
index \(i\) can satisfy this condition.
It is therefore one common index for the entire neighborhood.
By convexity of the ball and the assumed hull identity,
\begin{equation}\label{eq:one-ball-neighborhood}
 \partial P(y)\subset\clB{v_i}{r}
 \qquad(\abs{y-x}\le la).
\end{equation}

\emph{Transfer to the physical pressure and its transform.}
Take \(v\in\partial q(x)\) and a unit vector \(z\in\R^2\).
At \(y=x+la z\), choose any support of \(P\).
Using \eqref{eq:one-ball-neighborhood} to bound this support and
the approximation at \(x,y\), we get
\[
 la\,v\cdot z
 \le q(y)-q(x)
 \le P(y)-P(x)+2\delta
 \le la\,(v_i\cdot z+r)+2\delta.
\]
Testing every unit \(z\) proves
\[
 \abs{v-v_i}\le r+\frac{2\delta}{la}.
\]
By \Cref{k:subdifferential}, the generating supports of the active branch
of \(T_aq\) are obtained by adding \((au^2/2,0)\).
They all lie in \(\clB{v_i}{r^+}\) by
the definition of \(r^+\) and \(\abs{u}\le U_*\).
The exact convex-hull formula over all active branches gives
\eqref{eq:one-step-conclusion}.
\end{proof}

The common density band is what prevents supports near one active
shifted point from mixing different groups.  Once one group has been
selected, ordinary secants transfer its entire support bound across a
pressure error \(\delta\); division by the test distance \(la\)
explains the term \(2\delta/(la)\).  The remaining term in \(r^+\)
bounds the shift \((au^2/2,0)\) of each support under the transform.

\subsection{A common choice of scales}\label{m:schedule}

For \(R_1>2\), define
\begin{equation}\label{eq:mass-range}
 R_n=R_1^{2^{n-1}},\qquad
 a_1=A,\qquad a_n=R_n^{-1/16}\quad(n\ge2).
\end{equation}
Let
\begin{equation}\label{eq:truncated-models}
 \varphi_n=\psi_t-t\sum_{j=1}^n a_jw_{R_j},
 \qquad q_n=Q_{\varphi_n},
 \qquad P_n=T_{a_n}q_{n-1}\quad(n\ge1),
\end{equation}
where \(\varphi_0=\psi_t\) and \(q_0=q^{(t)}\).
Thus \(P_1=T_Aq^{(t)}\) is the pressure used in
\Cref{i:initial-separation}.

Choose \(R_1\) large enough initially that
\(\sum_{n\ge2}a_n<1\).
Every truncated potential then belongs to the stable family
\eqref{f:family}.
By \Cref{k:range}, for all sufficiently large \(R_1\),
\begin{equation}\label{eq:approximation-errors}
 \norm{q_n-P_n}_{W_1}\le\delta_n,
 \qquad \delta_n=C_3R_n^{-1/4}\quad(n\ge1),
\end{equation}
with one constant \(C_3\) independent of \(n\) and \(R_1\).
Indeed, the reference range is at most two for \(n=1\), giving the
stronger estimate \(O(R_1^{-1/2})\).
For \(n\ge2\), it is at most \(R_{n-1}=\sqrt{R_n}\).
The ratio threshold in \Cref{k:range} is uniform in the reference, so it
suffices to enlarge the same \(R_1\).

The crucial relations are
\begin{equation}\label{eq:summable-errors}
 a_{n+1}=R_n^{-1/8},\qquad
 \frac{\delta_n}{a_{n+1}}=C_3R_n^{-1/8},
 \qquad
 \sum_{n\ge1}\left(a_{n+1}+\delta_n+
             \frac{\delta_n}{a_{n+1}}\right)
       \longrightarrow0\quad(R_1\to\infty).
\end{equation}
For example, for every \(q>0\),
\[
 \sum_{n\ge1}R_n^{-q}
 =\sum_{n\ge1}(R_1^{-q})^{2^{n-1}}
 \le \frac{R_1^{-q}}{1-R_1^{-q}}.
\]
This explicit bound will meet all the conditions below.

By \Cref{k:transform-properties}, fix \(C_0>0\), independently of
\(R_1\), such that
\begin{equation}\label{eq:small-mass-increment}
 0\le T_aQ-Q\le C_0a\qquad\text{on }W_1
\end{equation}
for all family pressures and \(0<a\le A+1\).

Set
\[
 W_n=W\left(\frac14-S_*\sum_{j=1}^{n-1}a_{j+1}\right),
\]
and define
\begin{equation}\label{eq:slope-radii}
 r_1=\frac{\tau}{2},\qquad
 r_{n+1}=r_n+\frac{2\delta_n}{la_{n+1}}
                   +\frac12a_{n+1}U_*^2.
\end{equation}

We now enlarge \(R_1\) once more so that
\begin{align}
 S_*\sum_{n\ge1}a_{n+1}&<\frac18,
 &\sum_{n\ge1}\left(\frac{2\delta_n}{la_{n+1}}
                         +\frac12a_{n+1}U_*^2\right)&<\frac{\tau}{2},
 \label{eq:domain-radius-budget}\\
 \sup_{n\ge1}\frac{2\delta_n}{ka_{n+1}}&<\frac p{16},
 &\sum_{n\ge1}(\delta_n+C_0a_{n+1})&<\frac{\gamma}{64}.
 \label{eq:band-pressure-budget}
\end{align}
All four conditions follow from \eqref{eq:summable-errors};
no further choices of masses or ranges are made.
In particular, every \(W_n\) contains \(W(1/8)\) and \(r_n<\tau\).

The pressures stay uniformly close to the three-plane limit:
\begin{equation}\label{eq:pressure-near-limit}
 \norm{P_n-P_*}_{W_1}<\frac{\gamma}{32}\qquad(n\ge1).
\end{equation}
To see this, use
\[
 \norm{P_{n+1}-P_n}_{W_1}
 \le \norm{T_{a_{n+1}}q_n-q_n}_{W_1}
          +\norm{q_n-P_n}_{W_1}
 \le C_0a_{n+1}+\delta_n
\]
and telescope from the initial error \(\gamma/64\).
This comparison uses \eqref{eq:small-mass-increment} for the physical
pressure \(q_n\), so no shifted evaluation of the approximation
\eqref{eq:approximation-errors} is required here.

\begin{proposition}[Persistence at every scale]\label{m:propagation}
For every \(n\ge1\) and \(y\in W_n\),
\begin{equation}\label{eq:separated-generators}
 \partial P_n(y)=
 \conv\left(\partial P_n(y)\cap
                 \bigcup_{i=1}^3\clB{v_i}{r_n}\right).
\end{equation}
\end{proposition}

\begin{proof}
The case \(n=1\) is \Cref{i:initial-separation}.
Each \(q_n\) belongs to the family, and each \(P_n\) has the common
Lipschitz bound \(C_*\).  Apply \Cref{m:one-step} inductively with
\(P=P_n\), \(q=q_n\), \(a=a_{n+1}\), \(\delta=\delta_n\),
\(r=r_n\), and \(s\) equal to the radius of \(W_n\).
The first condition in \eqref{eq:domain-radius-budget} keeps the
parameter radius above \(1/8\), while the second gives \(r_n<\tau\).
The other hypotheses are \eqref{eq:approximation-errors} and the first
condition in \eqref{eq:band-pressure-budget}.  The new domain and
slope radius are exactly \(W_{n+1}\) and \(r_{n+1}\).
\end{proof}

The three groups now force a two-dimensional set of supports at one
parameter.  Minimizing against their barycenter first finds that
barycenter in a subdifferential; its separation from every two-group
hull then forces all three groups to occur there.

\begin{proposition}[A common ball of pressure supports]\label{m:pressure-ball}
For every \(n\ge1\), there is
\(\lambda_n\in\operatorname{int}W(1/8)\) such that
\begin{equation}\label{eq:common-pressure-ball}
 \clB{g}{\gamma}\subset\partial P_n(\lambda_n).
\end{equation}
\end{proposition}

\begin{proof}
Minimize \(P_n(\lambda)-g\cdot\lambda\) on the compact ball \(W(1/8)\).
For a boundary point, \Cref{i:geometry} and
\eqref{eq:pressure-near-limit} give
\begin{align*}
 &[P_n(\lambda)-g\cdot\lambda]
   -[P_n(\lambda_*)-g\cdot\lambda_*]\\
 &\hspace{2em}>
   \frac{\gamma}{8}-\frac{2\gamma}{32}
   =\frac{\gamma}{16}>0.
\end{align*}
Thus a minimizer \(\lambda_n\) is interior, and
\(g\in\partial P_n(\lambda_n)\).
By \Cref{m:propagation}, this subdifferential is generated by its
intersections with the three balls \(\clB{v_i}{r_n}\).
The vector \(g\) lies outside the convex hull of any two of the larger
balls \(\clB{v_i}{\tau}\).
Each of the three intersections must therefore be nonempty.
Choose one actual support from each.
Their triangle contains \(\clB{g}{\gamma}\) by \Cref{i:geometry}, and
their convex hull is contained in \(\partial P_n(\lambda_n)\).
\end{proof}

\subsection{The final potential and its decay}

Define one potential, now with all parameters fixed, by
\begin{equation}\label{eq:final-potential}
 \Phi(x)=\phi(\abs{x})
   =\psi_t(\abs{x})-t\sum_{n\ge1}a_nw_{R_n}(x).
\end{equation}

\begin{proposition}[Stability and algebraic decay]\label{m:tail}
The function \(\phi\) is bounded and continuous on \([0,\infty)\),
is stable with \(B=t^2/2\), and satisfies
\(\abs{\phi(r)}\le Cr^{-3-1/32}\) for all \(r\ge1\), where
\begin{equation}\label{eq:tail-constant}
 C=2^{3+1/32}t^2+t\norm{w}_{\infty}
       \left(A R_1^{1/32}
              +\sum_{n\ge2}R_1^{-2^{n-1}/32}\right)<\infty.
\end{equation}
\end{proposition}

\begin{proof}
Since \(\sum_na_n<A+1\) and \(R_n\ge1\), the attraction series
converges uniformly.
The core and every summand are bounded, continuous, and radial.
The common cell estimate of \Cref{f:stability} applies to the sum and
gives \(U^\Phi\ge-B_0N=-(t^2/2)N\).

Set \(\varepsilon=1/32\).
A nonzero summand at radius \(r\) must have \(R_n\ge r\), and then
\[
 r^{3+\varepsilon}R_n^{-3}\le R_n^\varepsilon.
\]
The core vanishes for \(r\ge1+1/t<2\) and is at most \(t^2\).
For every \(r\ge1\), therefore,
\[
 r^{3+\varepsilon}\abs{\phi(r)}
 \le 2^{3+\varepsilon}t^2
          +t\norm{w}_\infty\sum_{n:R_n\ge r}a_nR_n^\varepsilon
 \le C.
\]
Finally,
\[
 \sum_{n\ge1}a_nR_n^{1/32}
 =A R_1^{1/32}+\sum_{n\ge2}R_1^{-2^{n-1}/32}<\infty
\]
by the geometric-series bound following
\eqref{eq:summable-errors}.
\end{proof}

\subsection{From pressure supports to a temperature singularity}

By \Cref{f:thermodynamics}, the infinite potential has canonical and
grand-canonical thermodynamic limits.
By \Cref{f:tail-pressure} and \eqref{eq:approximation-errors},
\[
 \norm{P_n-Q_\Phi}_{W_1}
 \le \delta_n+
 C_{W_1,t}\sum_{j>n}a_j\longrightarrow0.
\]
Choose a subsequence of \(\lambda_n\) converging to
\(\lambda_c=(\beta_c,h_c)\in W(1/8)\).
In particular, \(\beta_c\in[7/8,9/8]\).
For each \(v\in\clB{g}{\gamma}\) and \(\lambda\in W_1\),
\eqref{eq:common-pressure-ball} gives
\[
 P_n(\lambda)\ge P_n(\lambda_n)+v\cdot(\lambda-\lambda_n).
\]
Uniform convergence and continuity pass this inequality to the limit.
It is a supporting inequality for \(Q_\Phi\) on a neighborhood of
\(\lambda_c\).
Such a local support extends globally on the convex domain \(\Omega\):
a violation at another point would, by convexity along the joining
segment, give a violation arbitrarily close to \(\lambda_c\).
Consequently,
\begin{equation}\label{eq:limiting-support-ball}
 \clB{g}{\gamma}\subset\partial Q_\Phi(\beta_c,h_c).
\end{equation}

We state explicitly the duality implication used to finish the proof.
It is the supporting-hyperplane form of Fenchel duality
\citep[Theorem~11.1 and Proposition~11.3]{RockafellarWets1998}, applied
only in the field coordinate and with the inverse temperature retained
as a parameter.

\begin{lemma}[From pressure supports to canonical slopes]\label{m:canonical-corner}
Suppose that \(H(\beta,d)\) and \(Q(\beta,h)\) have the duality and
finiteness properties of \Cref{f:thermodynamics}; assume also, as there,
that \(H(\,\cdot\,,d)\) is concave on \((0,\infty)\) for every \(d>0\).
If
\[
 (s_1,\rho),(s_2,\rho)\in\partial Q(\beta_c,h_c),
 \qquad \rho>0,\quad s_1<s_2,
\]
then \(H(\,\cdot\,,\rho)\) has finite one-sided derivatives at \(\beta_c\)
and
\[
 H'_+(\beta_c,\rho)\le-s_2<-s_1\le H'_-(\beta_c,\rho).
\]
\end{lemma}

\begin{proof}
The field part of either support inequality says that \(\rho\) supports
\(Q(\beta_c,\cdot)\) at \(h_c\).
Density duality therefore gives
\(H(\beta_c,\rho)=h_c\rho-Q(\beta_c,h_c)\).
The full support inequality, with \(s=s_i\), reads
\[
 Q(\beta,h)\ge Q(\beta_c,h_c)
                 +s(\beta-\beta_c)+\rho(h-h_c).
\]
Rearrange and take the supremum over \(h\).
The reverse duality gives
\[
 H(\beta,\rho)\le H(\beta_c,\rho)-s(\beta-\beta_c).
\]
The function \(H(\,\cdot\,,\rho)\) is finite and concave on
\((0,\infty)\), so its one-sided derivatives at an interior point
exist and are finite.
Dividing the last inequality by positive and negative increments yields
the stated ordering.
\end{proof}

\begin{proof}[Proof of \Cref{thm:main}]
Take the potential \eqref{eq:final-potential} and the constants in
\Cref{m:tail}.
Write \(g=(g_1,\rho_0)\), where \(\rho_0=5p/3>0\), and set
\[
 \eta=\min\{\gamma/2,\rho_0/2\},\qquad
 I=(\rho_0-\eta,\rho_0+\eta)\subset(0,\infty).
\]
For every \(\rho\in I\), the two points
\[
 (g_1-\gamma/2,\rho),\qquad(g_1+\gamma/2,\rho)
\]
lie in \(\clB{g}{\gamma}\), since their squared distances from \(g\)
are \((\gamma/2)^2+(\rho-\rho_0)^2<\gamma^2/2\).
By \eqref{eq:limiting-support-ball}, both support the same pressure at
the same \((\beta_c,h_c)\), independently of \(\rho\).
Lemma~\ref{m:canonical-corner} therefore gives finite one-sided
canonical slopes whose difference is at least \(\gamma\) for every
\(\rho\in I\).
The all-sequence canonical limit of \Cref{f:thermodynamics} applies
in particular to \(N_L=\lfloor\rho L^3\rfloor\) for every \(\beta>0\).
The specified free energy equals
\[
 f(\beta,\rho)=\frac{t}{\beta}H_\Phi(\beta,\rho).
\]
For each \(\rho\in I\), it is finite for every \(\beta>0\).
The finite concave function \(H_\Phi(\,\cdot\,,\rho)\) is continuous
at \(\beta_c\); hence the common derivative term from \(1/\beta\)
cancels in the difference, giving
\[
 f'_-(\beta_c,\rho)-f'_+(\beta_c,\rho)
 =\frac{t}{\beta_c}
       \bigl(H'_{\Phi,-}(\beta_c,\rho)
                    -H'_{\Phi,+}(\beta_c,\rho)\bigr)
 \ge\frac{t\gamma}{\beta_c}>0.
\]
The potential and the density interval are fixed before the limiting
critical parameter is selected, so this one \(\beta_c\) works throughout
\(I\).
\end{proof}


\bibliographystyle{plainnat}
\begingroup
\raggedright
\bibliography{refs}
\endgroup
\end{document}
